% This is samplepaper.tex, a sample chapter demonstrating the
% LLNCS macro package for Springer Computer Science proceedings;
% Version 2.20 of 2017/10/04
%\documentclass[runningheads]{llncs}
\documentclass[runningheads]{llncs}

\usepackage{amsmath}
\usepackage{booktabs}
\usepackage{siunitx}

\usepackage{tabularx}
\usepackage{makecell}
\usepackage{tikz}
\usepackage{adjustbox}
\newcommand{\Kheader}{\textbf{Header}} % Replace 'Header' with your actual header text
\usepackage{graphicx}
\usepackage{xcolor, colortbl}
\usepackage{placeins}
\usepackage{multirow}
\usepackage{array}
\usepackage{needspace}
\usepackage{orcidlink}
\hypersetup{hypertexnames=false}
\usepackage{longtable}
\usepackage{algorithm}
\usepackage{float}
\usepackage{algpseudocode}

\usepackage[T5]{fontenc}
\usepackage[utf8]{inputenc}
\newenvironment{red}{\color{red}}{}
\newenvironment{blue}{\color{blue}}{}
\usepackage{pgfplots}
\usepackage{subcaption}
\usepgfplotslibrary{fillbetween}
\usepgfplotslibrary{groupplots}
\pgfplotsset{compat=1.18}

% ── Per-architecture colors required by sensitivity_figures.tex ──────────
\definecolor{colorViT}{HTML}{D62728}
\definecolor{colorDenseNet}{HTML}{1F77B4}
\definecolor{colorResNet}{HTML}{2CA02C}
\definecolor{colorEfficientNet}{HTML}{FF7F0E}
\definecolor{colorMobileNet}{HTML}{9467BD}
\definecolor{colorVGG}{HTML}{7F7F7F}

\definecolor{hdrblue}{RGB}{45, 62, 80}      % Xanh đá phiến đậm (Slate) - sang và có chiều sâu
\definecolor{hdrtext}{RGB}{255, 255, 255}  % Chữ trắng
\definecolor{stripe}{RGB}{245, 247, 250}   % Màu sọc xám xanh cực nhẹ, đủ để phân biệt dòng

\usepackage{geometry}
\geometry{a4paper, margin=1in}

% -------------------------------------------------------------------------
\definecolor{rowalt}{RGB}{247,249,252}
\newcolumntype{K}{>{\centering\arraybackslash\small}p{1.55cm}}
\newcommand{\ms}[2]{\small$#1{\scriptstyle\,\pm#2}$}
\newcommand{\msb}[2]{\small$\mathbf{#1}{\scriptstyle\,\pm#2}$}
\newcommand{\bestms}[2]{\cellcolor{black!10}\msb{#1}{#2}}
\begin{document}
\raggedbottom
\title{Improving Image Classification with Validation-Guided Post-Training Checkpoint Averaging}

\author{
  % Minh-Dung Le\orcidlink{0009-0006-6924-4363} \and
  % Minh-Duc Hoang \and
  % Hoang-Vu Truong \and
  % Thi-Thu-Hong Phan\thanks{Corresponding author: hongptt@fpt.edu.vn}\orcidlink{0000-0001-6880-3721}
}
%
% \authorrunning{Le et al.}
\institute{
  % AIT laboratory, Faculty of Artificial Intelligence, \\ FPT University, Da Nang 550000, Viet Nam\\
}

\maketitle              % typeset the header of the

\begin{abstract}
  Selecting a single checkpoint based on validation loss can make the final
  model sensitive to fluctuations along the training trajectory. We propose a
  lightweight post-training strategy that ranks epoch checkpoints by validation
  loss, selects the top-$k$ candidates, and uniformly averages their parameters.
  The method requires no changes to the training pipeline and retains
  single-model inference cost. We evaluate the proposed strategy across six
  architectures and five image-classification datasets using five random seeds.
  At the common $k=5$ setting, it achieves higher mean accuracy and macro-F1
  than the single-checkpoint baseline in all 30 model--dataset configurations.
  When the configurations are weighted equally, mean accuracy increases from
  81.73\% to 83.76\%, while mean macro-F1 increases from 81.65\% to 83.86\%,
  corresponding to gains of 2.04 and 2.21 percentage points, respectively.
  Although the optimal aggregation size remains configuration-dependent, the
  results show that averaging a small validation-ranked checkpoint set can
  improve final performance and frequently reduce run-to-run variability
  without increasing inference-time complexity.
\end{abstract}

\section{Introduction}
Convolutional neural networks (CNNs) have achieved strong performance across image classification benchmarks, supported by continual advances in architecture design and optimization \cite{krizhevsky2012imagenet,simonyan2015vgg,he2016deep}. Despite these advances, most practical training pipelines still finalize a model via single-checkpoint selection: from a sequence of epoch-wise checkpoints, the state with the best validation accuracy or lowest validation loss is retained \cite{bishop2006pattern,prechelt2002early}. While simple and widely adopted, this protocol assumes that late-epoch training is sufficiently stable such that a single best epoch reliably reflects generalization performance.
In practice, late-epoch validation metrics often fluctuate due to mini-batch sampling noise, stochastic optimization dynamics, and data variability \cite{prechelt2002early,izmailov2019averagingweightsleadswider,garipov2018losssurfacesmodeconnectivity}. Checkpoints with similar training loss can exhibit noticeably different generalization behavior \cite{keskar2017largebatchtrainingdeeplearning}, making single-checkpoint selection sensitive to transient validation noise. A natural remedy is to aggregate information from multiple model states rather than committing to one checkpoint.
Prediction-time ensembling achieves this but increases inference cost proportionally to ensemble size \cite{dietterich2000ensemble}. Weight averaging—which averages parameters rather than outputs—preserves single-model inference. Foundational work on iterate averaging \cite{polyak1992acceleration} and Stochastic Weight Averaging (SWA) \cite{izmailov2019averagingweightsleadswider} has shown that averaging high-quality solutions encountered during training yields more stable models.
% SWA, however, requires explicit integration into the training procedure (specialized learning-rate schedules and online weight accumulation), making it inapplicable to retrospective post-training selection from standard runs.
SWA~\cite{izmailov2019averagingweightsleadswider} performs weight
averaging as part of an additional SGD phase governed by either a cyclical or
constant learning-rate schedule. Model states are sampled according to that
schedule and accumulated during optimization. In contrast, our method does not
introduce any additional optimization phase or prescribe a learning-rate regime;
it retrospectively ranks checkpoints from a completed standard training run
using validation loss.
This paper proposes a criterion-guided post-training checkpoint aggregation strategy that operates strictly at the model selection stage, requiring no changes to the optimization pipeline. We identify the top-$k$ checkpoints with the lowest validation loss and average their weights to form the final model. This approach extends standard validation-based selection from a winner-takes-all decision to aggregation over multiple strong candidates, reducing sensitivity to noisy late-stage fluctuations without increasing inference-time cost.
We evaluate the method on six widely used architectures: ViT-Base~\cite{dosovitskiy2021imageworth16x16words}, VGG16~\cite{simonyan2015vgg}, ResNet101~\cite{he2016deep}, MobileNetV2~\cite{sandler2018mobilenetv2}, DenseNet121~\cite{huang2017densely}, and EfficientNetB0~\cite{tan2019efficientnet} across five datasets: three agricultural image-classification datasets (Tomato Leaf Disease~\cite{solapure2024tomato}, Burmese Grape Leaf Disease~\cite{rahman2025burmese}, and Potato Leaf Disease~\cite{Shabrina2023Potato}), CIFAR-100~\cite{krizhevsky2009learning}, and Tiny ImageNet~\cite{le2015tiny}. Each model--dataset configuration is repeated using five random seeds, and
performance is reported using accuracy and macro-F1. This expanded evaluation examines both run-to-run stability and
generalization beyond the original agricultural domain.

The paper is organised as follows. Section~\ref{sec:related_work} reviews related work.
Section~\ref{sec:methodology} presents the proposed method.
Section~\ref{sec:data_description} describes the datasets.
Section~\ref{sec:ExpandResults} reports the experimental results.
Section~\ref{sec:discussion} discusses the findings, limitations, and implications of the proposed approach.
Finally, Section~\ref{sec:conclusion} concludes the paper and outlines future work.

\section{Related Work}
\label{sec:related_work}

% ----------------------------------------------------------
\subsection{Training-Time Weight Averaging}
% ----------------------------------------------------------

Early optimization theory established that averaging iterates reduces
variance and improves asymptotic convergence~\cite{polyak1992acceleration}.
In deep learning, Exponential Moving Average (EMA) applies this idea
online, maintaining a smoothed parameter copy throughout
training~\cite{tarvainen2018meanteachersbetterrole}.
Although effective, EMA is influenced by the entire training history,
potentially retaining information from earlier suboptimal states.

SWA~\cite{izmailov2019averagingweightsleadswider} averages model states
sampled along an SGD trajectory generated under either a cyclical or
constant learning-rate schedule.
Starting from a conventionally trained model, SWA continues optimization
and captures checkpoints at the minimum learning rate of each cycle, or
at every epoch under a constant learning rate.
Its canonical algorithm therefore prescribes learning-rate bounds, a cycle
length, and a fixed number of optimization iterations.
The original SWA study allows the preceding conventional training stage to
use either the full training budget or a predetermined reduced budget;
the latter corresponds to stopping conventional training before the full
budget and then starting the SWA phase, and should not be interpreted as
validation-based early stopping within SWA itself.

Latest Weight Averaging (LAWA)~\cite{kaddour2022stop} proposes maintaining
a sliding window of the $k$ most recent epoch checkpoints and replacing the
current model with their uniform average at regular intervals during
ongoing training.
On ResNet-50/ImageNet and RoBERTa/WikiText-103, this procedure advances
loss and accuracy curves by tens of epochs, yielding substantial GPU-hour
savings.
LAWA selects candidates by recency: the $k$ checkpoints are always the
latest ones, sampled at a fixed interval, regardless of their individual
validation performance.
Our method instead ranks all epoch checkpoints by validation loss and
selects the top-$k$ regardless of when they occur in the training
trajectory; a strong checkpoint from an earlier epoch is included while a
weaker recent one is not.

Trainable Weight Averaging (TWA)~\cite{li2022trainable} further extends the
direction by replacing pre-set uniform coefficients with learned ones.
TWA optimizes the averaging weights within the subspace spanned by the
candidate checkpoints, using either training or validation data
(TWA-v), and consistently improves over SWA.
The key distinction from our work is that TWA introduces a secondary
optimization stage to determine the aggregation coefficients, whereas our
method uses a uniform average over validation-loss-ranked checkpoints. It
requires no secondary optimization, although parameter averaging and Batch
Normalization recalibration introduce a small one-time post-training cost.

% ----------------------------------------------------------
\subsection{Inference-Time Ensembling}
% ----------------------------------------------------------

Classic model ensembling improves generalization by combining diverse
hypotheses, at the cost of inference overhead proportional to the number
of ensemble members~\cite{dietterich2000ensemble}.
Snapshot ensembles partially reduce training cost by exploiting
learning-rate schedules to obtain multiple diverse checkpoints within a
single run~\cite{huang2017snapshotensemblestrain1}.
In resource-constrained deployments, inference-time aggregation is often
undesirable due to its additional compute and memory requirements.

% ----------------------------------------------------------
\subsection{Post-Training Checkpoint Aggregation}
% ----------------------------------------------------------

Temporal-window averaging is a simple post-training strategy in which
checkpoints are selected according to their position in the training
sequence---for example, the final $k$ epochs or checkpoints sampled at
fixed time intervals.
This practice is well established in neural machine translation; notably,
Vaswani et al.~\cite{vaswani2017attention} report averaging the last 5
(base) or last 20 (big) checkpoints to obtain their final Transformer
models, and subsequent NMT systems have routinely adopted this
convention~\cite{junczys2016neural}.
The strategy assumes that temporally adjacent or recent checkpoints have
reached a stable region and are therefore similarly suitable for
aggregation.

However, late-stage validation performance is not necessarily monotonic.
A predefined temporal window may include checkpoints with degraded
validation performance caused by transient optimization fluctuations,
overfitting, or scheduler-induced changes.
At the same time, it may exclude stronger checkpoints located outside the
selected interval.
Consequently, the quality of the aggregated model depends on the training
endpoint, the window location, and the chosen interval rather than on the
observed generalization quality of the individual checkpoints.

Model Soups~\cite{wortsman2022model} revisit post-training weight
aggregation in the context of fine-tuning large pre-trained models.
Rather than selecting the single best checkpoint from a hyperparameter
sweep, the greedy soup variant sequentially adds each fine-tuned model to
the soup only if doing so improves validation accuracy.
Model Soups therefore also uses a validation-guided criterion to filter
candidates before averaging.
The key structural difference from our method is that Model Soups
operates \emph{across independently trained models} with different
hyperparameter configurations, whereas our approach operates
\emph{within a single training run}, selecting from the epoch checkpoints
produced by that run according to their validation loss.

% ----------------------------------------------------------
\subsection{Positioning of the Present Work}
% ----------------------------------------------------------

The proposed method differs from SWA not only in when averaging is
performed but also in how candidate checkpoints are generated and
selected.
Canonical SWA deliberately produces averaging candidates through an
additional optimization phase using a cyclical or constant learning-rate
schedule.
LAWA integrates checkpoint averaging as an online procedure inside the
training loop.
TWA introduces a learnable weighting stage that optimizes aggregation
coefficients on held-out data.
Model Soups operates across independent runs with varied hyperparameters.

Our method instead operates after the original training run has terminated
and ranks the epoch checkpoints generated by that run according to
validation loss.
Consequently, the proposed strategy does not prescribe an optimizer,
learning-rate schedule, cycle structure, stopping rule, or secondary
optimization.
It is compatible with standard training pipelines using different
schedulers and with or without early stopping, provided that epoch
checkpoints and their validation losses are available.
Unlike temporal-window averaging, which selects checkpoints by epoch
order, our approach averages the top-$k$ checkpoints with the lowest
validation losses, which need not be consecutive or restricted to the
final epochs.

\FloatBarrier

\section{Methodology}
\label{sec:methodology}

\subsection{Overview}

Figure~\ref{fig:proposed_method} illustrates the proposed post-training checkpoint aggregation framework. The proposed framework operates strictly at the model selection stage of a standard image classification training pipeline. All epoch checkpoints and their validation losses are saved during training without affecting the optimization process. After training terminates, a validation-loss-guided criterion is applied to select and aggregate multiple high-performing checkpoints into a single set of averaged parameters. No component of the training process, including the optimizer, learning-rate schedule, early stopping, or training duration, is modified. The resulting model requires no additional inference cost compared to a single-checkpoint baseline.
\FloatBarrier

\begin{figure}[H]
  \centering
  \includegraphics[width=0.8\linewidth]{Pipeline.png}
  \caption{Workflow of the proposed post-training checkpoint aggregation strategy. The baseline selects a single checkpoint with the lowest validation loss, whereas the proposed method averages the top-$k$ selected checkpoints and recomputes Batch Normalization statistics when applicable before final evaluation.}
  \label{fig:proposed_method}
\end{figure}

\FloatBarrier
The detailed formulation of the proposed criterion-guided checkpoint aggregation strategy is presented in the following subsection.

\subsection{Proposed Post-Training Checkpoint Aggregation}

Let
$\{\mathbf{w}_1,\mathbf{w}_2,\ldots,\mathbf{w}_T\}$
denote the trainable-parameter vectors saved at the end of each training
epoch, where $T$ is the total number of completed epochs. Let
$\ell_t^{\mathrm{val}}$ denote the loss of checkpoint $\mathbf{w}_t$
on the held-out checkpoint-selection set.

The conventional baseline selects the single checkpoint with the lowest
validation loss:
\begin{equation}
  t^*
  =
  \arg\min_{1\leq t\leq T}\ell_t^{\mathrm{val}},
  \qquad
  \mathbf{w}_{\mathrm{baseline}}
  =
  \mathbf{w}_{t^*}.
\end{equation}

The proposed method instead identifies the indices of the $k$ checkpoints
with the lowest validation losses:
\begin{equation}
  \mathcal{S}_k
  =
  \operatorname{TopKLowestIndices}
  \left(
    \{\ell_t^{\mathrm{val}}\}_{t=1}^{T},k
  \right).
\end{equation}

Their trainable parameters are then averaged element-wise:
\begin{equation}
  \mathbf{w}_{\mathrm{avg}}
  =
  \frac{1}{|\mathcal{S}_k|}
  \sum_{t\in\mathcal{S}_k}\mathbf{w}_t.
\end{equation}

Here, $\mathcal{S}_k\subseteq\{1,\ldots,T\}$ contains the epoch indices
corresponding to the $k$ smallest validation losses, with ties broken by
preferring later epochs. The averaging is applied only to trainable
parameters. For networks using Batch Normalization
(BN)~\cite{ioffe2015batch}, checkpoint-specific running statistics may not
correspond to the averaged parameters. Following
SWA~\cite{izmailov2019averagingweightsleadswider}, the BN running statistics
are therefore reset and recomputed from the training data after averaging.
This step requires only a forward pass, without gradients or optimizer
updates, and the finalized model retains single-model inference cost.

% \FloatBarrier
% \subsection{Evaluated Architectures}
% To assess the generality of the proposed strategy, we evaluate six architectures spanning convolutional and transformer-based design paradigms. All pretrained backbone parameters remain trainable and are fine-tuned
% end-to-end on each target dataset. The baseline and all Top-$k$ variants use identical architectures, initializations, and training trajectories, so that their only difference is the post-training checkpoint-selection rule.

% ViT-Base~\cite{dosovitskiy2021imageworth16x16words} reformulates image classification as a sequence modeling problem by splitting an input image into fixed-size patches processed by a Transformer encoder, enabling capture of long-range spatial dependencies. VGG16~\cite{simonyan2015vgg} provides a strong reference baseline via uniform $3\times3$ convolutional layers followed by fully connected classifiers. ResNet101~\cite{he2016deep} addresses deep network optimization through residual connections that facilitate gradient propagation. DenseNet121~\cite{huang2017densely} employs dense connectivity in which each layer receives feature maps from all preceding layers, promoting feature reuse and parameter efficiency. MobileNetV2~\cite{sandler2018mobilenetv2} uses depthwise separable convolutions to achieve competitive accuracy under tight computational budgets. EfficientNetB0~\cite{tan2019efficientnet} introduces compound scaling of depth, width, and resolution via fixed scaling coefficients, achieving strong accuracy--efficiency trade-offs.

\subsection{Evaluated Architectures}
To assess the generality of the proposed strategy, we evaluate six architectures spanning convolutional and transformer-based paradigms. All pretrained backbone parameters are fine-tuned end-to-end. The baseline and all Top-$k$ variants use identical architectures, initializations, and training trajectories, differing solely in their post-training checkpoint-selection rule.

We evaluate: (i) \textbf{ViT-Base}~\cite{dosovitskiy2021imageworth16x16words} as the representative transformer encoder for capturing long-range spatial dependencies; (ii) \textbf{VGG16}~\cite{simonyan2015vgg} as a standard baseline with uniform convolutions; (iii) \textbf{ResNet101}~\cite{he2016deep} utilizing residual connections; (iv) \textbf{DenseNet121}~\cite{huang2017densely} employing dense connectivity for feature reuse; (v) \textbf{MobileNetV2}~\cite{sandler2018mobilenetv2} featuring depthwise separable convolutions for tight computational budgets; and (vi) \textbf{EfficientNetB0}~\cite{tan2019efficientnet} leveraging compound scaling for optimal accuracy--efficiency trade-offs.

\subsection{Top-$k$ Checkpoint Aggregation Procedure}

Algorithm~\ref{alg:checkpoint_aggregation} summarizes the proposed
post-training procedure.

\begin{algorithm}[ht]
  \caption{Post-Training Top-$k$ Validation-Loss-Guided Checkpoint Aggregation}
  \label{alg:checkpoint_aggregation}
  \begin{algorithmic}[1]

    \Require Completed checkpoint set
    $\mathcal{C}
    =
    \{(\ell_t^{\mathrm{val}},\mathbf{w}_t)\}_{t=1}^{T}$,
    aggregation size $k$, training data $\mathcal{D}_{\mathrm{train}}$

    \Ensure Final aggregated model $\mathcal{M}_{\mathrm{avg}}$

    \State $\mathcal{S}_k \leftarrow
    \operatorname{TopKLowestIndices}
    \left(
      \{\ell_t^{\mathrm{val}}\}_{t=1}^{T},k
    \right)$
    \Comment{ties favor later epochs}

    \State $\mathbf{w}_{\mathrm{avg}} \leftarrow
    \dfrac{1}{|\mathcal{S}_k|}
    \displaystyle\sum_{t\in\mathcal{S}_k}\mathbf{w}_t$

    \State Load $\mathbf{w}_{\mathrm{avg}}$ into
    $\mathcal{M}_{\mathrm{avg}}$

    \If{$\mathcal{M}_{\mathrm{avg}}$ contains BN layers}
    \State Reset the BN running statistics
    \State Recompute them using one forward pass over
    $\mathcal{D}_{\mathrm{train}}$ in training mode
    \Statex \hspace{\algorithmicindent}
    without gradients or optimizer updates
    \EndIf

    \State \Return $\mathcal{M}_{\mathrm{avg}}$

  \end{algorithmic}
\end{algorithm}
\FloatBarrier
\section{Datasets and Evaluation Protocols}
\label{sec:data_description}

We evaluate the proposed method on five image-classification datasets:
three agricultural datasets and two general-purpose benchmarks.
Table~\ref{tab:dataset_summary} summarizes their principal characteristics
and evaluation protocols.

\begin{table}[!ht]
  \centering
  \caption{Summary of datasets and evaluation protocols. ``Selection'' denotes
    the split used for checkpoint ranking, and ``Final evaluation'' the held-out
    split on which the reported metrics are computed. Tiny ImageNet has no
    publicly labelled test split, so its official validation set serves as the
    held-out evaluation set and its results are reported as validation metrics
  throughout~\cite{abai2019densenet,huynh2022vision}.}
  \label{tab:dataset_summary}
  \resizebox{\linewidth}{!}{
    \begin{tabular}{lccccc}
      \toprule
      Dataset & Classes & Available images & Train & Selection & Final evaluation \\
      \midrule
      Tomato Leaf Disease
      & 10 & 1,609
      & 70\% & 15\% & 15\% test \\

      Burmese Grape Leaf Disease
      & 5 & 3,103
      & 70\% & 15\% & 15\% test \\

      Potato Leaf Disease
      & 7 & 3,076
      & 70\% & 15\% & 15\% test \\

      CIFAR-100
      & 100 & 50,000 train + 10,000 test
      & 45,000 & 5,000 & 10,000 official test \\

      Tiny ImageNet
      & 200 & 100,000 train + 10,000 official val.
      & 90,000 & 10,000 & 10,000 official val.$^{\dagger}$ \\
      \bottomrule
      \multicolumn{6}{@{}l}{\footnotesize $^{\dagger}$Reported as validation
      metrics; test labels are not publicly available.}
      \\
  \end{tabular}}
\end{table}

\subsection{Agricultural Image Datasets}
The Tomato~\cite{solapure2024tomato}, Burmese Grape~\cite{rahman2025burmese}, and Potato~\cite{Shabrina2023Potato} Leaf Disease datasets contain 1,609 images across ten classes, 3,103 images across five classes, and 3,076 images across seven classes, respectively. All images are resized to a standardized resolution fitting the input requirements of each evaluated architecture. These datasets provide a diverse agricultural benchmark to evaluate the robustness of the proposed checkpoint-selection strategy under practical, domain-specific scenarios.
\subsection{CIFAR-100}

CIFAR-100~\cite{krizhevsky2009learning} contains 50,000 training images and
10,000 official test images distributed evenly across 100 classes. The original
training partition is divided using a stratified split into 45,000 training
images and 5,000 checkpoint-selection images. The official test set is kept
intact and used only for final evaluation.

\subsection{Tiny ImageNet}

Tiny ImageNet~\cite{le2015tiny} comprises 100,000 training and 10,000 official
validation images across 200 classes. The labels of the official test split are
not publicly released, so the official validation set takes the role of the
held-out evaluation set, following the protocol adopted by prior work on this
benchmark~\cite{abai2019densenet,huynh2022vision}. Because that split is a
validation split by construction, every number we report for Tiny ImageNet is
labelled as a validation metric (validation accuracy, validation macro-F1)
rather than a test metric, and it is never used for checkpoint ranking.
Checkpoint selection instead uses a separate split obtained by stratifying the
original training partition into 90,000 local-training images and 10,000
local-validation images, identically across all experimental setups. The two
splits are therefore disjoint, and the evaluation split remains unseen by the
selection criterion exactly as for the other four datasets.
\FloatBarrier

% \section{Experiments and Results}
% \label{sec:ExpandResults}
% \subsection{Training Configuration}
% All models are initialized with ImageNet-pretrained weights and fine-tuned end-to-end with randomly initialized classifier heads. Experiments are conducted on an NVIDIA Tesla T4 GPU using the Kaggle platform. The training pipeline—including optimizer, learning-rate schedule, early stopping, and total epochs—is kept identical between the baseline and the proposed method, as detailed in Tables~\ref{tab:tomato_configs_transposed}--\ref{tab:potato_configs_transposed}. All models use Linear Warmup ($0.06 \times$ total epochs) followed by Cosine Annealing ($\eta_{\min}=1 \times 10^{-6}$). Checkpoints and their validation losses are recorded at the end of each epoch; the aggregation is applied strictly post-training.

% Each model dataset configuration is trained independently using five random
% seeds: 1, 10, 42, 100, and 500. For each seed, the single-checkpoint baseline
% and all Top-$k$ variants are derived from the same training trajectory and
% checkpoint set. Therefore, the compared methods differ only in their
% post-training model-selection rule. Results are reported as mean $\pm$
% standard deviation across the five runs.

\section{Experiments and Results}
\label{sec:ExpandResults}

\subsection{Training Configuration}

All architectures are instantiated using the \texttt{timm}
library~\cite{rw2019timm}, from which the corresponding pretrained
weights are obtained. The training pipeline---including the optimizer, learning-rate schedule, early stopping, and total epochs---is kept identical between the baseline and the proposed method. All models utilize a Linear Warmup followed by Cosine Annealing. Checkpoints and their corresponding validation losses are recorded at the end of each epoch, and the aggregation is applied strictly post-training.

Settings that are shared by every model--dataset configuration are stated once
here rather than repeated per row. All models are optimized with
Adam and a cross-entropy objective with label smoothing $0.1$, at an input
resolution of $224\times224$. The learning-rate schedule is a linear warmup
over the first $6\%$ of the scheduled epochs followed by cosine annealing to
$\eta_{\min}=1\times10^{-6}$. Each backbone is fitted with the same classifier
head---two hidden layers of 256 and 128 units with dropout $0.5$---and weight
decay is $1\times10^{-6}$ for all models except ViT-Base, which uses
$1\times10^{-7}$. Table~\ref{tab:training_configs} therefore lists only the
four quantities that actually vary across configurations: batch size, epoch
budget, peak learning rate, and early-stopping patience. Reported epoch
budgets are the scheduled maxima; early stopping frequently terminates a run
before that budget is reached, and the number of checkpoints actually produced
is reported in Section~\ref{sec:selection_geometry}.

Each model--dataset configuration is trained independently using five random
seeds: 1, 10, 42, 100, and 500. For each seed, the single-checkpoint baseline
and all Top-$k$ variants are derived from the same training trajectory and
completed checkpoint set. Therefore, they differ only in the post-training
checkpoint-selection rule. Results are reported as mean $\pm$ standard
deviation across the five runs. Unless stated otherwise, ``final evaluation''
refers to the official test split for CIFAR-100 and the held-out 15\% test
split for the three agricultural datasets; for Tiny ImageNet it refers to the
official validation split, and the corresponding numbers are reported as
validation metrics (Section~\ref{sec:data_description}).
{%
  \small
  \setlength{\tabcolsep}{5pt} % Giảm nhẹ một chút để fit chuẩn form canh giữa
  \renewcommand{\arraystretch}{1.2}

  % Định nghĩa lại toàn bộ cột: Căn giữa theo cả chiều dọc (m) và chiều ngang (centering)
  \newcolumntype{D}{>{\centering\arraybackslash}m{2.5cm}}
  \newcolumntype{M}{>{\centering\arraybackslash}m{2.7cm}}

  % Định nghĩa các cột số phía sau cũng dùng kiểu 'm' căn giữa đồng bộ
  \newcolumntype{B}{>{\centering\arraybackslash}m{1.2cm}} % Batch
  \newcolumntype{E}{>{\centering\arraybackslash}m{1.2cm}} % Epochs
  \newcolumntype{L}{>{\centering\arraybackslash}m{1.5cm}} % LR
  \newcolumntype{P}{>{\centering\arraybackslash}m{1.4cm}} % Patience

  \begin{longtable}{@{} D M B E L P @{}}

    \caption{Model-specific training configurations across all evaluated datasets. \textit{Patience} denotes the early-stopping epochs.}
    \label{tab:training_configs}\\[-4pt]

    %% First-page header
    \toprule
    \rowcolor{rowalt}
    \textbf{Dataset} &
    \textbf{Model}   &
    \textbf{Batch}   &
    \textbf{Epochs}  &
    \textbf{LR}      &
    \textbf{Patience}\\[3pt]
    \midrule
    \endfirsthead

    %% Continuation header
    \multicolumn{6}{@{}l}{%
      \footnotesize\itshape
    \tablename~\thetable{} --- continued from previous page}\\[4pt]
    \toprule
    \rowcolor{rowalt}
    \textbf{Dataset} &
    \textbf{Model}   &
    \textbf{Batch}   &
    \textbf{Epochs}  &
    \textbf{LR}      &
    \textbf{Patience}\\[3pt]
    \midrule
    \endhead

    %% Continuation footer
    \midrule
    \multicolumn{6}{r@{}}{%
    \footnotesize\itshape Continued on next page\ldots}\\
    \endfoot

    %% Last footer
    \bottomrule
    \endlastfoot

    % TOMATO LEAF
    & ViT-Base       & 32  & 50  & $9{\times}10^{-7}$   & 10 \\*
    \rowcolor{stripe}
    & VGG16          & 64  & 150 & $5{\times}10^{-4}$   & 15 \\*
    & ResNet101      & 64  & 150 & $5{\times}10^{-4}$   & 15 \\*
    \rowcolor{stripe}
    & MobileNetV2    & 32  & 150 & $1.5{\times}10^{-4}$ & 15 \\*
    & DenseNet121    & 64  & 100 & $1.5{\times}10^{-5}$ & 10 \\*
    \rowcolor{stripe}
    \multirow{-6}{*}{\textbf{Tomato Leaf}}
    & EfficientNetB0 & 64  & 150 & $5{\times}10^{-4}$   & 15 \\
    \midrule

    % BURMESE GRAPE
    & ViT-Base       & 64  & 50  & $2{\times}10^{-5}$   & 10 \\*
    \rowcolor{stripe}
    & VGG16          & 16  & 100 & $2{\times}10^{-5}$   & 15 \\*
    & ResNet101      & 64  & 50  & $2{\times}10^{-5}$   & 10 \\*
    \rowcolor{stripe}
    & MobileNetV2    & 64  & 100 & $1{\times}10^{-3}$   & 20 \\*
    & DenseNet121    & 64  & 100 & $1{\times}10^{-5}$   & 15 \\*
    \rowcolor{stripe}
    \multirow{-6}{*}{\makecell{\textbf{Burmese}\\\textbf{Grape}}}
    & EfficientNetB0 & 64  & 100 & $2{\times}10^{-3}$   & 15 \\
    \midrule

    % POTATO LEAF
    & ViT-Base       & 64  & 50  & $2{\times}10^{-5}$   & 10 \\*
    \rowcolor{stripe}
    & VGG16          & 32  & 70  & $5{\times}10^{-4}$   & 10 \\*
    & ResNet101      & 64  & 50  & $6{\times}10^{-4}$   & 10 \\*
    \rowcolor{stripe}
    & MobileNetV2    & 32  & 150 & $1.5{\times}10^{-4}$ & 15 \\*
    & DenseNet121    & 32  & 50  & $5{\times}10^{-4}$   & 10 \\*
    \rowcolor{stripe}
    \multirow{-6}{*}{\textbf{Potato Leaf}}
    & EfficientNetB0 & 64  & 50  & $6{\times}10^{-4}$   & 10 \\
    \midrule

    % CIFAR-100
    & ViT-Base       & 64  & 40  & $2{\times}10^{-5}$   &  8 \\*
    \rowcolor{stripe}
    & VGG16          & 64  & 100 & $5{\times}10^{-5}$   & 12 \\*
    & ResNet101      & 128 & 60  & $3{\times}10^{-4}$   & 12 \\*
    \rowcolor{stripe}
    & MobileNetV2    & 128 & 60  & $3{\times}10^{-4}$   & 12 \\*
    & DenseNet121    & 128 & 60  & $3{\times}10^{-4}$   & 12 \\*
    \rowcolor{stripe}
    \multirow{-6}{*}{\textbf{CIFAR-100}}
    & EfficientNetB0 & 128 & 60  & $5{\times}10^{-4}$   & 12 \\
    \midrule

    % TINY IMAGENET
    & ViT-Base       & 1024 & 60 & $2{\times}10^{-4}$   & 10 \\*
    \rowcolor{stripe}
    & VGG16          & 1024 & 50 & $1{\times}10^{-4}$   & 10 \\*
    & ResNet101      & 1024 & 70 & $2{\times}10^{-4}$   & 10 \\*
    \rowcolor{stripe}
    & MobileNetV2    & 1024 & 60 & $2{\times}10^{-4}$   & 10 \\*
    & DenseNet121    & 1024 & 60 & $2{\times}10^{-4}$   & 10 \\*
    \rowcolor{stripe}
    \multirow{-6}{*}{\textbf{Tiny ImageNet}}
    & EfficientNetB0 & 1024 & 60 & $2{\times}10^{-4}$   & 10 \\

  \end{longtable}
}%

%

\subsection{Performance of the Proposed Strategy}
\label{sec:main_results}
Tables~\ref{tab:results-cifar100}--\ref{tab:results-potato} report
final-evaluation accuracy and macro-F1 for six architectures across five
datasets; for Tiny ImageNet these are validation metrics, as explained in
Section~\ref{sec:data_description}. At least one Top-$k$ setting improves mean
accuracy over the single-checkpoint baseline for every model--dataset
configuration. Furthermore, the common $k=5$ setting improves both mean
accuracy and macro-F1 across all 30 configurations. Macro-precision,
macro-recall, and AUC were also recorded; because they follow the same trend
as accuracy and macro-F1 in every configuration, they are summarised in
aggregate in Table~\ref{tab:aggregate_k_results} rather than repeated in the
per-dataset tables.

\paragraph{General trend.}
The effect of $k$ is not uniformly monotonic. While many configurations
continue to improve through $k=4$ or $k=5$, others saturate earlier or
exhibit small decreases after an intermediate optimum. The highest observed
accuracy occurs at $k=4$ or $k=5$ in 25 of the 30 configurations, suggesting
that these values provide a useful practical range rather than a universally
optimal choice.

\paragraph{General-purpose benchmarks.}
On the two general-purpose benchmarks, CIFAR-100 (Table~\ref{tab:results-cifar100}), ViT-Base achieves the
strongest absolute performance, reaching $89.64\pm0.28\%$ accuracy at $k{=}5$
versus $87.19\pm0.65\%$ at baseline—an improvement of 2.45 percentage points
accompanied by a reduction in variance.
EfficientNetB0 gains 4.39 percentage points ($76.76 \to 81.15\%$), the largest absolute
accuracy improvement among the CNNs on this dataset.
On Tiny ImageNet (Table~\ref{tab:results-tiny-imagenet}), whose entries are
validation metrics because the benchmark releases no labelled test split,
ResNet101 exhibits the most pronounced gain: $77.50 \to 81.34\%$ validation
accuracy (+3.84 percentage points at $k{=}5$),
while ViT-Base advances from $89.68$ to $91.53\%$ with markedly reduced
standard deviation ($0.34 \to 0.12$), indicating greater run-to-run stability.

\paragraph{Agricultural datasets.}
On the three domain-specific benchmarks
(Tables~\ref{tab:results-burmese-grape}--\ref{tab:results-potato}), the
magnitude of the gain varies more substantially across architectures, and
the run-to-run standard deviations are generally larger than on the two
general-purpose benchmarks.
On Tomato Leaf, ResNet101 shows the largest improvement: $81.84 \to 87.52\%$
accuracy (+5.68 percentage points) and $84.48 \to 90.01\%$ macro-F1 (+5.53 percentage points), both at $k{=}4$.
On Potato Leaf, DenseNet121 improves from $86.79$ to $90.26\%$ (+3.47 percentage points),
and ViT-Base reaches $90.90\%$ accuracy at $k{=}4$ with the lowest
standard deviation among the six model-specific peak-accuracy entries
($\pm0.49$).
Burmese Grape yields smaller and more architecture-dependent gains;
MobileNetV2 nonetheless
improves from $83.89$ to $87.19\%$ (+3.30 percentage points) at $k{=}4$.

\paragraph{Run-to-run variability.}
Checkpoint aggregation frequently, but not universally, reduces run-to-run
variability. At $k=5$, the accuracy standard deviation is lower than that of
the baseline in 25 of the 30 configurations, and the macro-F1 standard
deviation is lower in 23 configurations. Averaged over all 30 configurations,
the across-seed standard deviation falls from $1.43$ to $1.11$ percentage
points for accuracy ($-22.3\%$) and from $1.45$ to $1.11$ for macro-F1
($-23.5\%$); the reduction is largest on the two general-purpose benchmarks
($-46\%$ on CIFAR-100 and $-56\%$ on Tiny ImageNet) and smallest on Tomato
Leaf ($-8\%$), the dataset with the fewest images and the noisiest baseline.
Notable individual reductions include
Potato--ViT-Base ($1.51\rightarrow0.40$), Burmese Grape--VGG16
($2.04\rightarrow0.98$), and CIFAR-100--ResNet101
($1.54\rightarrow0.53$). Several configurations exhibit small increases,
so variance reduction should be interpreted as a general empirical tendency
rather than a guaranteed property.
% ── CIFAR-100 ────────────────────────────────────────────────────
\begin{table}[htbp]
  \centering
  \caption{Results on CIFAR-100 (mean $\pm$ standard deviation).
  Bold gray cells indicate the highest mean.}
  \label{tab:results-cifar100}
  \setlength{\tabcolsep}{3pt}
  \renewcommand{\arraystretch}{1.16}
  \scriptsize
  \begin{tabular}{@{}ll KKKKKK@{}}
    \toprule
    \textbf{Metric} & \textbf{$k$}
    & \textbf{VGG16}
    & \makecell{\textbf{Dense}\\\textbf{Net121}}
    & \makecell{\textbf{Efficient}\\\textbf{NetB0}}
    & \makecell{\textbf{Mobile}\\\textbf{NetV2}}
    & \makecell{\textbf{ResNet}\\\textbf{101}}
    & \makecell{\textbf{ViT}\\\textbf{Base}} \\
    \midrule

    \multirow{5}{*}{\textbf{Accuracy}}
    & 1 & \ms{71.57}{0.41} & \ms{75.39}{0.20} & \ms{76.76}{0.51}
    & \ms{73.23}{0.41} & \ms{75.84}{1.54} & \ms{87.19}{0.65} \\
    & 2 & \ms{72.10}{0.45} & \ms{78.10}{0.20} & \ms{79.85}{0.31}
    & \ms{74.24}{0.22} & \ms{78.08}{0.28} & \ms{88.68}{0.28} \\
    & 3 & \ms{72.34}{0.40} & \ms{78.73}{0.12} & \ms{80.51}{0.06}
    & \ms{74.68}{0.24} & \ms{78.60}{0.36} & \ms{89.12}{0.59} \\
    & 4 & \bestms{72.48}{0.47} & \ms{78.93}{0.18} & \ms{80.92}{0.22}
    & \ms{74.92}{0.18} & \bestms{78.82}{0.48} & \ms{89.39}{0.36} \\
    & 5 & \ms{72.46}{0.43} & \bestms{79.23}{0.30} & \bestms{81.15}{0.31}
    & \bestms{75.06}{0.15} & \ms{78.80}{0.53} & \bestms{89.64}{0.28} \\

    \midrule

    \multirow{5}{*}{\textbf{Macro-F1}}
    & 1 & \ms{71.59}{0.42} & \ms{75.41}{0.28} & \ms{76.75}{0.52}
    & \ms{73.25}{0.50} & \ms{75.92}{1.52} & \ms{87.01}{0.69} \\
    & 2 & \ms{72.11}{0.46} & \ms{78.09}{0.22} & \ms{79.82}{0.31}
    & \ms{74.27}{0.22} & \ms{78.15}{0.26} & \ms{88.52}{0.32} \\
    & 3 & \ms{72.36}{0.42} & \ms{78.72}{0.16} & \ms{80.46}{0.08}
    & \ms{74.71}{0.24} & \ms{78.67}{0.35} & \ms{88.95}{0.65} \\
    & 4 & \bestms{72.50}{0.48} & \ms{78.92}{0.19} & \ms{80.88}{0.25}
    & \ms{74.95}{0.20} & \bestms{78.87}{0.43} & \ms{89.25}{0.39} \\
    & 5 & \ms{72.48}{0.45} & \bestms{79.21}{0.28} & \bestms{81.10}{0.34}
    & \bestms{75.04}{0.18} & \bestms{78.87}{0.50} & \bestms{89.52}{0.30} \\
    \bottomrule
  \end{tabular}
\end{table}

% ── Tiny ImageNet ─────────────────────────────────────────────────
\begin{table}[htbp]
  \centering
  \caption{Results on Tiny ImageNet (mean $\pm$ standard deviation). Because
    the test labels of this benchmark are not publicly available, the official
    validation split is used as the held-out evaluation set and all entries are
    validation metrics~\cite{abai2019densenet,huynh2022vision}. Bold gray cells
  indicate the highest mean.}
  \label{tab:results-tiny-imagenet}
  \setlength{\tabcolsep}{3pt}
  \renewcommand{\arraystretch}{1.16}
  \scriptsize
  \begin{tabular}{@{}ll KKKKKK@{}}
    \toprule
    \textbf{Metric} & \textbf{$k$}
    & \textbf{VGG16}
    & \makecell{\textbf{Dense}\\\textbf{Net121}}
    & \makecell{\textbf{Efficient}\\\textbf{NetB0}}
    & \makecell{\textbf{Mobile}\\\textbf{NetV2}}
    & \makecell{\textbf{ResNet}\\\textbf{101}}
    & \makecell{\textbf{ViT}\\\textbf{Base}} \\
    \midrule

    \multirow{5}{*}{\makecell{\textbf{Validation}\\\textbf{Accuracy}}}
    & 1 & \ms{67.54}{0.85} & \ms{76.82}{0.61} & \ms{80.32}{0.32}
    & \ms{72.90}{0.54} & \ms{77.50}{0.60} & \ms{89.68}{0.34} \\
    & 2 & \ms{68.04}{0.46} & \ms{77.74}{0.11} & \ms{80.64}{0.25}
    & \ms{73.39}{0.38} & \ms{80.41}{0.10} & \ms{90.96}{0.15} \\
    & 3 & \ms{68.12}{0.43} & \ms{77.88}{0.13} & \ms{80.72}{0.26}
    & \ms{73.40}{0.31} & \ms{80.86}{0.13} & \ms{91.31}{0.10} \\
    & 4 & \bestms{68.35}{0.38} & \bestms{78.06}{0.15} & \bestms{80.81}{0.19}
    & \ms{73.46}{0.23} & \ms{81.23}{0.14} & \ms{91.43}{0.10} \\
    & 5 & \ms{68.34}{0.35} & \ms{78.04}{0.22} & \ms{80.78}{0.20}
    & \bestms{73.53}{0.21} & \bestms{81.34}{0.32} & \bestms{91.53}{0.12} \\

    \midrule

    \multirow{5}{*}{\makecell{\textbf{Validation}\\\textbf{Macro-F1}}}
    & 1 & \ms{67.36}{0.87} & \ms{76.81}{0.59} & \ms{80.25}{0.32}
    & \ms{72.74}{0.54} & \ms{77.58}{0.51} & \ms{89.69}{0.33} \\
    & 2 & \ms{67.85}{0.48} & \ms{77.73}{0.10} & \ms{80.58}{0.26}
    & \ms{73.24}{0.38} & \ms{80.41}{0.10} & \ms{90.96}{0.15} \\
    & 3 & \ms{67.91}{0.46} & \ms{77.85}{0.13} & \ms{80.66}{0.27}
    & \ms{73.25}{0.32} & \ms{80.85}{0.16} & \ms{91.31}{0.10} \\
    & 4 & \bestms{68.14}{0.42} & \bestms{78.03}{0.16} & \bestms{80.75}{0.19}
    & \ms{73.31}{0.23} & \ms{81.23}{0.15} & \ms{91.43}{0.10} \\
    & 5 & \ms{68.13}{0.39} & \ms{78.01}{0.21} & \ms{80.72}{0.21}
    & \bestms{73.38}{0.21} & \bestms{81.34}{0.34} & \bestms{91.52}{0.12} \\
    \bottomrule
  \end{tabular}
\end{table}

% ── Burmese Grape ────────────────────────────────────────────────
\begin{table}[htbp]
  \centering
  \caption{Results on Burmese Grape (mean $\pm$ standard deviation).
  Bold gray cells indicate the highest mean.}
  \label{tab:results-burmese-grape}
  \setlength{\tabcolsep}{3pt}
  \renewcommand{\arraystretch}{1.16}
  \scriptsize
  \begin{tabular}{@{}ll KKKKKK@{}}
    \toprule
    \textbf{Metric} & \textbf{$k$}
    & \textbf{VGG16}
    & \makecell{\textbf{Dense}\\\textbf{Net121}}
    & \makecell{\textbf{Efficient}\\\textbf{NetB0}}
    & \makecell{\textbf{Mobile}\\\textbf{NetV2}}
    & \makecell{\textbf{ResNet}\\\textbf{101}}
    & \makecell{\textbf{ViT}\\\textbf{Base}} \\
    \midrule

    \multirow{5}{*}{\textbf{Accuracy}}
    & 1 & \ms{83.76}{2.04} & \ms{84.62}{1.62} & \ms{82.74}{2.09}
    & \ms{83.89}{1.54} & \ms{85.38}{0.94} & \ms{86.98}{1.92} \\
    & 2 & \ms{85.30}{1.39} & \ms{84.96}{1.34} & \ms{85.56}{1.27}
    & \ms{85.94}{2.28} & \ms{85.13}{1.49} & \ms{86.58}{2.18} \\
    & 3 & \ms{85.43}{0.84} & \ms{85.00}{1.50} & \ms{85.77}{0.74}
    & \ms{86.67}{1.53} & \ms{85.17}{1.61} & \bestms{87.05}{1.44} \\
    & 4 & \ms{85.77}{0.93} & \ms{84.96}{1.40} & \bestms{86.54}{1.05}
    & \bestms{87.19}{1.54} & \bestms{85.68}{1.41} & \ms{86.79}{1.47} \\
    & 5 & \bestms{85.94}{0.98} & \bestms{85.17}{1.50} & \ms{86.45}{1.15}
    & \ms{87.01}{1.67} & \ms{85.60}{1.28} & \ms{87.01}{1.50} \\

    \midrule

    \multirow{5}{*}{\textbf{Macro-F1}}
    & 1 & \ms{81.83}{3.42} & \ms{82.44}{1.46} & \ms{81.17}{2.70}
    & \ms{81.85}{1.62} & \ms{84.21}{1.17} & \ms{84.85}{1.49} \\
    & 2 & \ms{83.97}{1.68} & \ms{83.01}{1.25} & \ms{84.39}{1.65}
    & \ms{84.55}{1.79} & \ms{84.02}{1.82} & \ms{85.70}{2.57} \\
    & 3 & \ms{83.69}{1.80} & \ms{83.01}{1.32} & \ms{84.84}{0.70}
    & \ms{85.42}{1.18} & \ms{84.08}{1.24} & \bestms{86.00}{1.87} \\
    & 4 & \ms{84.08}{1.67} & \ms{82.84}{1.29} & \ms{85.49}{0.84}
    & \ms{85.64}{1.24} & \bestms{84.54}{1.24} & \ms{85.66}{1.75} \\
    & 5 & \bestms{84.41}{1.42} & \bestms{83.06}{1.39} & \bestms{85.51}{0.90}
    & \bestms{86.08}{1.39} & \ms{84.40}{1.07} & \ms{85.89}{2.17} \\
    \bottomrule
  \end{tabular}
\end{table}

% ── Tomato Leaf ──────────────────────────────────────────────────
\begin{table}[htbp]
  \centering
  \caption{Results on Tomato Leaf (mean $\pm$ standard deviation).
  Bold gray cells indicate the highest mean.}
  \label{tab:results-tomato}
  \setlength{\tabcolsep}{3pt}
  \renewcommand{\arraystretch}{1.16}
  \scriptsize
  \begin{tabular}{@{}ll KKKKKK@{}}
    \toprule
    \textbf{Metric} & \textbf{$k$}
    & \textbf{VGG16}
    & \makecell{\textbf{Dense}\\\textbf{Net121}}
    & \makecell{\textbf{Efficient}\\\textbf{NetB0}}
    & \makecell{\textbf{Mobile}\\\textbf{NetV2}}
    & \makecell{\textbf{ResNet}\\\textbf{101}}
    & \makecell{\textbf{ViT}\\\textbf{Base}} \\
    \midrule

    \multirow{5}{*}{\textbf{Accuracy}}
    & 1 & \ms{82.08}{2.79} & \ms{84.08}{2.50} & \ms{83.20}{2.46}
    & \ms{83.04}{3.92} & \ms{81.84}{2.62} & \ms{83.12}{2.93} \\
    & 2 & \ms{83.36}{2.46} & \bestms{84.56}{2.32} & \ms{84.96}{2.81}
    & \ms{84.40}{2.34} & \ms{85.36}{1.74} & \ms{83.52}{2.86} \\
    & 3 & \bestms{84.00}{2.56} & \ms{84.48}{2.66} & \ms{85.20}{2.96}
    & \ms{85.12}{2.45} & \ms{86.64}{1.94} & \bestms{83.92}{1.87} \\
    & 4 & \bestms{84.00}{2.10} & \ms{84.48}{2.53} & \ms{85.52}{2.55}
    & \ms{85.92}{3.44} & \bestms{87.52}{2.57} & \ms{83.68}{2.34} \\
    & 5 & \ms{83.20}{2.68} & \ms{84.40}{2.29} & \bestms{85.60}{2.30}
    & \bestms{86.64}{3.51} & \ms{87.12}{2.56} & \ms{83.68}{2.43} \\

    \midrule

    \multirow{5}{*}{\textbf{Macro-F1}}
    & 1 & \ms{84.71}{2.67} & \ms{87.52}{1.94} & \ms{86.19}{2.21}
    & \ms{86.78}{3.99} & \ms{84.48}{1.92} & \ms{86.71}{2.77} \\
    & 2 & \ms{86.04}{2.34} & \bestms{87.98}{1.88} & \ms{88.52}{2.42}
    & \ms{87.84}{2.22} & \ms{88.08}{0.57} & \ms{87.23}{2.83} \\
    & 3 & \ms{86.87}{2.23} & \ms{87.91}{2.25} & \ms{88.58}{2.42}
    & \ms{88.56}{1.75} & \ms{89.38}{0.95} & \bestms{87.37}{2.03} \\
    & 4 & \bestms{86.96}{1.91} & \ms{87.94}{2.21} & \ms{88.78}{2.24}
    & \ms{88.93}{2.52} & \bestms{90.01}{1.84} & \ms{87.21}{2.67} \\
    & 5 & \ms{86.43}{1.96} & \ms{87.95}{1.94} & \bestms{88.99}{1.70}
    & \bestms{89.70}{2.80} & \ms{89.91}{1.58} & \ms{87.26}{2.65} \\
    \bottomrule
  \end{tabular}
\end{table}

% ── Potato Leaf ──────────────────────────────────────────────────
\begin{table}[htbp]
  \centering
  \caption{Results on Potato Leaf (mean $\pm$ standard deviation).
  Bold gray cells indicate the highest mean.}
  \label{tab:results-potato}
  \setlength{\tabcolsep}{3pt}
  \renewcommand{\arraystretch}{1.16}
  \scriptsize
  \begin{tabular}{@{}ll KKKKKK@{}}
    \toprule
    \textbf{Metric} & \textbf{$k$}
    & \textbf{VGG16}
    & \makecell{\textbf{Dense}\\\textbf{Net121}}
    & \makecell{\textbf{Efficient}\\\textbf{NetB0}}
    & \makecell{\textbf{Mobile}\\\textbf{NetV2}}
    & \makecell{\textbf{ResNet}\\\textbf{101}}
    & \makecell{\textbf{ViT}\\\textbf{Base}} \\
    \midrule

    \multirow{5}{*}{\textbf{Accuracy}}
    & 1 & \ms{86.03}{1.67} & \ms{86.79}{1.88} & \ms{85.94}{0.76}
    & \ms{87.09}{1.67} & \ms{86.20}{1.02} & \ms{90.30}{1.51} \\
    & 2 & \ms{87.18}{1.08} & \ms{89.10}{1.65} & \ms{88.12}{1.33}
    & \ms{88.03}{0.80} & \ms{88.42}{1.14} & \ms{90.30}{0.82} \\
    & 3 & \bestms{88.08}{1.03} & \ms{89.66}{1.34} & \ms{88.46}{0.97}
    & \ms{88.29}{0.87} & \ms{88.89}{1.42} & \ms{90.43}{0.51} \\
    & 4 & \ms{87.61}{1.38} & \bestms{90.26}{2.01} & \ms{88.68}{0.74}
    & \ms{88.50}{1.22} & \ms{88.76}{0.96} & \bestms{90.90}{0.49} \\
    & 5 & \ms{87.78}{1.26} & \ms{90.00}{2.15} & \bestms{88.97}{0.69}
    & \bestms{88.76}{0.72} & \bestms{89.10}{0.79} & \ms{90.60}{0.40} \\

    \midrule

    \multirow{5}{*}{\textbf{Macro-F1}}
    & 1 & \ms{84.52}{2.20} & \ms{84.90}{1.90} & \ms{84.67}{0.94}
    & \ms{85.52}{1.94} & \ms{84.84}{0.68} & \ms{88.06}{1.31} \\
    & 2 & \ms{86.07}{1.09} & \ms{87.18}{1.61} & \ms{86.80}{1.85}
    & \ms{86.89}{0.87} & \ms{87.53}{1.92} & \ms{88.29}{1.19} \\
    & 3 & \bestms{87.19}{1.23} & \ms{87.64}{1.44} & \ms{87.13}{1.66}
    & \ms{86.92}{1.13} & \ms{87.23}{2.54} & \ms{88.53}{0.83} \\
    & 4 & \ms{86.68}{1.52} & \bestms{88.53}{2.22} & \ms{87.29}{1.29}
    & \ms{87.21}{1.19} & \ms{87.50}{1.70} & \bestms{89.15}{1.21} \\
    & 5 & \ms{86.88}{1.75} & \ms{88.38}{2.41} & \bestms{87.57}{1.04}
    & \bestms{87.51}{0.81} & \bestms{87.59}{1.55} & \ms{88.94}{1.14} \\
    \bottomrule
  \end{tabular}
\end{table}

\FloatBarrier
\subsection{Sensitivity to the Aggregation Size}
\label{sec:k_sensitivity}

The effect of the aggregation size is architecture- and dataset-dependent.
Across the 30 model--dataset configurations, the highest observed mean
accuracy is obtained at $k=5$ in 14 configurations, at $k=4$ in 11,
at $k=3$ in four, and at $k=2$ in one. Thus, $k\in\{4,5\}$ accounts for
25 of the 30 accuracy-maximizing configurations. A similar pattern is
observed for macro-F1: the highest mean is achieved at $k=5$ in 16
configurations, at $k=4$ in ten, at $k=3$ in three, and at $k=2$ in one.

Nevertheless, performance is not uniformly monotonic as $k$ increases.
Mean accuracy is non-decreasing from $k=1$ to $k=5$ in 12 of the 30
configurations, while macro-F1 is non-decreasing in 14 configurations.
Several models therefore saturate at $k=3$ or $k=4$ and exhibit small
decreases when an additional checkpoint is included. Examples include
VGG16 and ResNet101 on CIFAR-100; VGG16, DenseNet121, and EfficientNetB0
on Tiny ImageNet; and multiple architectures on the agricultural datasets.

Despite these configuration-specific variations, the common $k=5$
setting improves mean accuracy and macro-F1 over the single-checkpoint
baseline in all 30 configurations. Table~\ref{tab:aggregate_k_results}
summarizes the equally weighted mean performance across all evaluated
model--dataset configurations.

% Per-dataset sensitivity curves (accuracy and macro-F1 vs. k, all six
% architectures). These duplicate Tables~3--7 in graphical form; uncomment the
% line below to include all five figures, or copy individual figure blocks out
% of sensitivity_figures.tex if page budget is tight.
% % ══════════════════════════════════════════════════════════════════════════
%  SENSITIVITY ANALYSIS FIGURES — All 5 datasets × Accuracy + F1-Score
%  Drop this file into your main .tex and % ══════════════════════════════════════════════════════════════════════════
%  SENSITIVITY ANALYSIS FIGURES — All 5 datasets × Accuracy + F1-Score
%  Drop this file into your main .tex and % ══════════════════════════════════════════════════════════════════════════
%  SENSITIVITY ANALYSIS FIGURES — All 5 datasets × Accuracy + F1-Score
%  Drop this file into your main .tex and \input{sensitivity_figures}
%  Requires in preamble: pgfplots, pgfplotsset{compat=1.18},
%    fillbetween, subcaption, xcolor
% ══════════════════════════════════════════════════════════════════════════

% ── Colors (already in your preamble — kept here for reference) ──────────
% \definecolor{colorViT}{HTML}{D62728}
% \definecolor{colorDenseNet}{HTML}{1F77B4}
% \definecolor{colorResNet}{HTML}{2CA02C}
% \definecolor{colorEfficientNet}{HTML}{FF7F0E}
% \definecolor{colorMobileNet}{HTML}{9467BD}
% \definecolor{colorVGG}{HTML}{7F7F7F}

% ── Shared axis style ────────────────────────────────────────────────────
\pgfplotsset{
  sens plot/.style={
    width=\linewidth,
    height=6.2cm,
    font=\small\rmfamily,
    xmin=0.55, xmax=5.45,
    xtick={1,2,3,4,5},
    xmajorgrids=true, ymajorgrids=true,
    grid style={dashed, gray!25},
    axis line style={gray!60},
    title style={font=\small\bfseries, yshift=-2pt},
    xlabel={Value of $k$\enspace{\footnotesize($k{=}1$: baseline)}},
    xlabel style={font=\small, yshift=3pt},
    ylabel style={font=\small},
    tick label style={font=\footnotesize},
    legend style={
      at={(0.5,-0.20)}, anchor=north,
      legend columns=3,
      font=\fontsize{6.8}{8}\selectfont,
      draw=gray!30, rounded corners=2pt,
      inner sep=3pt,
      /tikz/every even column/.append style={column sep=4pt},
    },
    line width=1.5pt,
    mark size=2.4pt,
    clip=false,   % allow annotations outside axis box
  }
}

% ── Annotation macro ─────────────────────────────────────────────────────
% \Ann{anchor}{x}{y}{val}{std}{color}
\newcommand{\Ann}[6]{%
  \node[anchor=#1, text=#6,
        font=\fontsize{5.8}{7}\selectfont\bfseries,
        fill=white, fill opacity=0.88, text opacity=1,
        inner sep=1pt, rounded corners=1.2pt]
    at (axis cs:#2,#3)
    {$#4$\\[-1.5pt]\normalfont$\pm\!#5$};}

% ── Shared legend entry (used in every subplot) ───────────────────────────
\newcommand{\SensLegend}{%
  \addlegendentry{VGG16}
  \addlegendentry{MobileNetV2}
  \addlegendentry{DenseNet121}
  \addlegendentry{ResNet101}
  \addlegendentry{EfficientNetB0}
  \addlegendentry{ViT-Base}
}


% ════════════════════════════════════════════════════════════════════════
%  HELPER: one subplot axis
%  Args: #1=ylabel text  #2=ymin  #3=ymax  #4=axis body (addplot + Ann)
% ════════════════════════════════════════════════════════════════════════
% (We don't use a separate macro for the axis body since pgfplots
%  needs the content inline — we just repeat the axis skeleton below.)


% ╔══════════════════════════════════════════════════════════════════════╗
% ║  FIGURE 1 — CIFAR-100                                               ║
% ╚══════════════════════════════════════════════════════════════════════╝
\begin{figure}[htbp]
\centering
\begin{subfigure}[t]{0.49\textwidth}
  \begin{adjustbox}{max width=\linewidth}
  \begin{tikzpicture}
  \begin{axis}[sens plot,
    title={CIFAR-100 — Accuracy},
    ylabel={Accuracy (\%)},
    ymin=69.5, ymax=91.5]

    \addplot[colorVGG,   mark=x]            coordinates{(1,71.57)(2,72.10)(3,72.34)(4,72.48)(5,72.46)};
    \addplot[colorMobileNet, mark=triangle] coordinates{(1,73.23)(2,74.24)(3,74.68)(4,74.92)(5,75.06)};
    \addplot[colorDenseNet,  mark=square]   coordinates{(1,75.39)(2,78.10)(3,78.73)(4,78.93)(5,79.23)};
    \addplot[colorResNet,    mark=triangle, mark options={rotate=180}]
                                            coordinates{(1,75.84)(2,78.08)(3,78.60)(4,78.82)(5,78.80)};
    \addplot[colorEfficientNet, mark=diamond]
                                            coordinates{(1,76.76)(2,79.85)(3,80.51)(4,80.92)(5,81.15)};
    \addplot[colorViT,   mark=*]            coordinates{(1,87.19)(2,88.68)(3,89.12)(4,89.39)(5,89.64)};

    % NOTE: a partial "std band" overlay for VGG16 was removed here. It was the
    % only band in the whole file and \closedcycle fills down to the axis rather
    % than between the two curves, so it rendered as a shaded wedge under VGG16
    % alone. The per-point std is already given by the \Ann annotations.

    % ── Annotations: baseline k=1 & peak ──
    \Ann{south}{1}{70.9}{71.57}{0.41}{colorVGG}
    \Ann{north}{1}{73.6}{73.23}{0.41}{colorMobileNet}
    \Ann{north}{5}{75.4}{75.06}{0.15}{colorMobileNet}
    \Ann{south}{1}{74.7}{75.39}{0.20}{colorDenseNet}
    \Ann{north}{5}{79.6}{79.23}{0.30}{colorDenseNet}
    \Ann{north}{1}{76.4}{75.84}{1.54}{colorResNet}
    \Ann{north}{1}{77.2}{76.76}{0.51}{colorEfficientNet}
    \Ann{north}{5}{81.6}{81.15}{0.31}{colorEfficientNet}
    \Ann{north}{1}{87.7}{87.19}{0.65}{colorViT}
    \Ann{north}{5}{90.1}{89.64}{0.28}{colorViT}

    \SensLegend
  \end{axis}
  \end{tikzpicture}
  \end{adjustbox}
  \caption{Test Accuracy vs.\ $k$}
  \label{fig:cifar100_acc}
\end{subfigure}\hfill
\begin{subfigure}[t]{0.49\textwidth}
  \begin{adjustbox}{max width=\linewidth}
  \begin{tikzpicture}
  \begin{axis}[sens plot,
    title={CIFAR-100 — F1-Score},
    ylabel={F1-Score (\%)},
    ymin=69.5, ymax=91.5]

    \addplot[colorVGG,         mark=x]            coordinates{(1,71.59)(2,72.11)(3,72.36)(4,72.50)(5,72.48)};
    \addplot[colorMobileNet,   mark=triangle]     coordinates{(1,73.25)(2,74.27)(3,74.71)(4,74.95)(5,75.04)};
    \addplot[colorDenseNet,    mark=square]       coordinates{(1,75.41)(2,78.09)(3,78.72)(4,78.92)(5,79.21)};
    \addplot[colorResNet,      mark=triangle, mark options={rotate=180}]
                                                  coordinates{(1,75.92)(2,78.15)(3,78.67)(4,78.87)(5,78.87)};
    \addplot[colorEfficientNet, mark=diamond]     coordinates{(1,76.75)(2,79.82)(3,80.46)(4,80.88)(5,81.10)};
    \addplot[colorViT,         mark=*]            coordinates{(1,87.01)(2,88.52)(3,88.95)(4,89.25)(5,89.52)};

    \Ann{south}{1}{70.9}{71.59}{0.42}{colorVGG}
    \Ann{north}{1}{73.7}{73.25}{0.50}{colorMobileNet}
    \Ann{north}{5}{75.4}{75.04}{0.18}{colorMobileNet}
    \Ann{south}{1}{74.7}{75.41}{0.28}{colorDenseNet}
    \Ann{north}{5}{79.6}{79.21}{0.28}{colorDenseNet}
    \Ann{north}{1}{76.5}{75.92}{1.52}{colorResNet}
    \Ann{north}{1}{77.2}{76.75}{0.52}{colorEfficientNet}
    \Ann{north}{5}{81.6}{81.10}{0.34}{colorEfficientNet}
    \Ann{north}{1}{87.6}{87.01}{0.69}{colorViT}
    \Ann{north}{5}{90.0}{89.52}{0.30}{colorViT}

    \SensLegend
  \end{axis}
  \end{tikzpicture}
  \end{adjustbox}
  \caption{Test F1-Score vs.\ $k$}
  \label{fig:cifar100_f1}
\end{subfigure}
\caption{Sensitivity analysis on \textbf{CIFAR-100}: effect of checkpoint-averaging
  parameter $k$ on Test Accuracy and F1-Score. $k{=}1$ is the single-checkpoint
  baseline. Annotated values show mean\,$\pm$\,std over five independent runs.}
\label{fig:sensitivity_cifar100}
\end{figure}

% ╔══════════════════════════════════════════════════════════════════════╗
% ║  FIGURE 2 — TINY IMAGENET                                           ║
% ╚══════════════════════════════════════════════════════════════════════╝
\begin{figure}[htbp]
\centering
\begin{subfigure}[t]{0.49\textwidth}
  \begin{adjustbox}{max width=\linewidth}
  \begin{tikzpicture}
  \begin{axis}[sens plot,
    title={Tiny ImageNet — Validation Accuracy},
    ylabel={Accuracy (\%)},
    ymin=65.5, ymax=93.0]

    \addplot[colorVGG,         mark=x]        coordinates{(1,67.54)(2,68.04)(3,68.12)(4,68.35)(5,68.34)};
    \addplot[colorMobileNet,   mark=triangle] coordinates{(1,72.90)(2,73.39)(3,73.40)(4,73.46)(5,73.53)};
    \addplot[colorDenseNet,    mark=square]   coordinates{(1,76.82)(2,77.74)(3,77.88)(4,78.06)(5,78.04)};
    \addplot[colorResNet,      mark=triangle, mark options={rotate=180}]
                                              coordinates{(1,77.50)(2,80.41)(3,80.86)(4,81.23)(5,81.34)};
    \addplot[colorEfficientNet,mark=diamond]  coordinates{(1,80.32)(2,80.64)(3,80.72)(4,80.81)(5,80.78)};
    \addplot[colorViT,         mark=*]        coordinates{(1,89.68)(2,90.96)(3,91.31)(4,91.43)(5,91.53)};

    \Ann{south}{1}{66.5}{67.54}{0.85}{colorVGG}
    \Ann{south}{1}{71.7}{72.90}{0.54}{colorMobileNet}
    \Ann{north}{5}{73.9}{73.53}{0.21}{colorMobileNet}
    \Ann{south}{1}{75.8}{76.82}{0.61}{colorDenseNet}
    \Ann{north}{5}{78.5}{78.04}{0.22}{colorDenseNet}
    \Ann{south}{1}{76.4}{77.50}{0.60}{colorResNet}
    \Ann{north}{5}{81.8}{81.34}{0.32}{colorResNet}
    \Ann{north}{1}{80.9}{80.32}{0.32}{colorEfficientNet}
    \Ann{north}{1}{90.2}{89.68}{0.34}{colorViT}
    \Ann{north}{5}{92.0}{91.53}{0.12}{colorViT}

    \SensLegend
  \end{axis}
  \end{tikzpicture}
  \end{adjustbox}
  \caption{Validation Accuracy vs.\ $k$}
  \label{fig:tinyimagenet_acc}
\end{subfigure}\hfill
\begin{subfigure}[t]{0.49\textwidth}
  \begin{adjustbox}{max width=\linewidth}
  \begin{tikzpicture}
  \begin{axis}[sens plot,
    title={Tiny ImageNet — Validation Macro-F1},
    ylabel={F1-Score (\%)},
    ymin=65.5, ymax=93.0]

    \addplot[colorVGG,         mark=x]        coordinates{(1,67.36)(2,67.85)(3,67.91)(4,68.14)(5,68.13)};
    \addplot[colorMobileNet,   mark=triangle] coordinates{(1,72.74)(2,73.24)(3,73.25)(4,73.31)(5,73.38)};
    \addplot[colorDenseNet,    mark=square]   coordinates{(1,76.81)(2,77.73)(3,77.85)(4,78.03)(5,78.01)};
    \addplot[colorResNet,      mark=triangle, mark options={rotate=180}]
                                              coordinates{(1,77.58)(2,80.41)(3,80.85)(4,81.23)(5,81.34)};
    \addplot[colorEfficientNet,mark=diamond]  coordinates{(1,80.25)(2,80.58)(3,80.66)(4,80.75)(5,80.72)};
    \addplot[colorViT,         mark=*]        coordinates{(1,89.69)(2,90.96)(3,91.31)(4,91.43)(5,91.52)};

    \Ann{south}{1}{66.3}{67.36}{0.87}{colorVGG}
    \Ann{south}{1}{71.6}{72.74}{0.54}{colorMobileNet}
    \Ann{north}{5}{73.8}{73.38}{0.21}{colorMobileNet}
    \Ann{south}{1}{75.7}{76.81}{0.59}{colorDenseNet}
    \Ann{north}{5}{78.4}{78.01}{0.21}{colorDenseNet}
    \Ann{south}{1}{76.4}{77.58}{0.51}{colorResNet}
    \Ann{north}{5}{81.8}{81.34}{0.34}{colorResNet}
    \Ann{north}{1}{80.8}{80.25}{0.32}{colorEfficientNet}
    \Ann{north}{1}{90.2}{89.69}{0.33}{colorViT}
    \Ann{north}{5}{92.0}{91.52}{0.12}{colorViT}

    \SensLegend
  \end{axis}
  \end{tikzpicture}
  \end{adjustbox}
  \caption{Validation Macro-F1 vs.\ $k$}
  \label{fig:tinyimagenet_f1}
\end{subfigure}
\caption{Sensitivity analysis on \textbf{Tiny ImageNet}. The official validation split is used as the held-out evaluation set because test labels are unavailable, so both panels report validation metrics.}
\label{fig:sensitivity_tinyimagenet}
\end{figure}

% ╔══════════════════════════════════════════════════════════════════════╗
% ║  FIGURE 3 — BURMESE GRAPE                                           ║
% ╚══════════════════════════════════════════════════════════════════════╝
\begin{figure}[htbp]
\centering
\begin{subfigure}[t]{0.49\textwidth}
  \begin{adjustbox}{max width=\linewidth}
  \begin{tikzpicture}
  \begin{axis}[sens plot,
    title={Burmese Grape — Accuracy},
    ylabel={Accuracy (\%)},
    ymin=79.5, ymax=90.5]

    \addplot[colorVGG,         mark=x]        coordinates{(1,83.76)(2,85.30)(3,85.43)(4,85.77)(5,85.94)};
    \addplot[colorMobileNet,   mark=triangle] coordinates{(1,83.89)(2,85.94)(3,86.67)(4,87.19)(5,87.01)};
    \addplot[colorDenseNet,    mark=square]   coordinates{(1,84.62)(2,84.96)(3,85.00)(4,84.96)(5,85.17)};
    \addplot[colorResNet,      mark=triangle, mark options={rotate=180}]
                                              coordinates{(1,85.38)(2,85.13)(3,85.17)(4,85.68)(5,85.60)};
    \addplot[colorEfficientNet,mark=diamond]  coordinates{(1,82.74)(2,85.56)(3,85.77)(4,86.54)(5,86.45)};
    \addplot[colorViT,         mark=*]        coordinates{(1,86.98)(2,86.58)(3,87.05)(4,86.79)(5,87.01)};

    \Ann{south}{1}{82.6}{83.76}{2.04}{colorVGG}
    \Ann{north}{5}{86.3}{85.94}{0.98}{colorVGG}
    \Ann{south}{1}{82.7}{83.89}{1.54}{colorMobileNet}
    \Ann{north}{4}{87.6}{87.19}{1.54}{colorMobileNet}
    \Ann{north}{1}{85.0}{84.62}{1.62}{colorDenseNet}
    \Ann{south}{1}{84.3}{85.38}{0.94}{colorResNet}
    \Ann{south}{1}{81.5}{82.74}{2.09}{colorEfficientNet}
    \Ann{north}{4}{87.0}{86.54}{1.05}{colorEfficientNet}
    \Ann{north}{1}{87.5}{86.98}{1.92}{colorViT}

    \SensLegend
  \end{axis}
  \end{tikzpicture}
  \end{adjustbox}
  \caption{Test Accuracy vs.\ $k$}
  \label{fig:grape_acc}
\end{subfigure}\hfill
\begin{subfigure}[t]{0.49\textwidth}
  \begin{adjustbox}{max width=\linewidth}
  \begin{tikzpicture}
  \begin{axis}[sens plot,
    title={Burmese Grape — F1-Score},
    ylabel={F1-Score (\%)},
    ymin=78.0, ymax=90.0]

    \addplot[colorVGG,         mark=x]        coordinates{(1,81.83)(2,83.97)(3,83.69)(4,84.08)(5,84.41)};
    \addplot[colorMobileNet,   mark=triangle] coordinates{(1,81.85)(2,84.55)(3,85.42)(4,85.64)(5,86.08)};
    \addplot[colorDenseNet,    mark=square]   coordinates{(1,82.44)(2,83.01)(3,83.01)(4,82.84)(5,83.06)};
    \addplot[colorResNet,      mark=triangle, mark options={rotate=180}]
                                              coordinates{(1,84.21)(2,84.02)(3,84.08)(4,84.54)(5,84.40)};
    \addplot[colorEfficientNet,mark=diamond]  coordinates{(1,81.17)(2,84.39)(3,84.84)(4,85.49)(5,85.51)};
    \addplot[colorViT,         mark=*]        coordinates{(1,84.85)(2,85.70)(3,86.00)(4,85.66)(5,85.89)};

    \Ann{south}{1}{80.6}{81.83}{3.42}{colorVGG}
    \Ann{north}{5}{84.8}{84.41}{1.42}{colorVGG}
    \Ann{south}{1}{80.6}{81.85}{1.62}{colorMobileNet}
    \Ann{north}{5}{86.5}{86.08}{1.39}{colorMobileNet}
    \Ann{north}{1}{82.9}{82.44}{1.46}{colorDenseNet}
    \Ann{north}{1}{84.7}{84.21}{1.17}{colorResNet}
    \Ann{south}{1}{79.9}{81.17}{2.70}{colorEfficientNet}
    \Ann{north}{5}{85.9}{85.51}{0.90}{colorEfficientNet}
    \Ann{north}{1}{85.3}{84.85}{1.49}{colorViT}
    \Ann{north}{3}{86.5}{86.00}{1.87}{colorViT}

    \SensLegend
  \end{axis}
  \end{tikzpicture}
  \end{adjustbox}
  \caption{Test F1-Score vs.\ $k$}
  \label{fig:grape_f1}
\end{subfigure}
\caption{Sensitivity analysis on \textbf{Burmese Grape}.}
\label{fig:sensitivity_grape}
\end{figure}

% ╔══════════════════════════════════════════════════════════════════════╗
% ║  FIGURE 4 — TOMATO                                                  ║
% ╚══════════════════════════════════════════════════════════════════════╝
\begin{figure}[htbp]
\centering
\begin{subfigure}[t]{0.49\textwidth}
  \begin{adjustbox}{max width=\linewidth}
  \begin{tikzpicture}
  \begin{axis}[sens plot,
    title={Tomato — Accuracy},
    ylabel={Accuracy (\%)},
    ymin=78.5, ymax=91.0]

    \addplot[colorVGG,         mark=x]        coordinates{(1,82.08)(2,83.36)(3,84.00)(4,84.00)(5,83.20)};
    \addplot[colorMobileNet,   mark=triangle] coordinates{(1,83.04)(2,84.40)(3,85.12)(4,85.92)(5,86.64)};
    \addplot[colorDenseNet,    mark=square]   coordinates{(1,84.08)(2,84.56)(3,84.48)(4,84.48)(5,84.40)};
    \addplot[colorResNet,      mark=triangle, mark options={rotate=180}]
                                              coordinates{(1,81.84)(2,85.36)(3,86.64)(4,87.52)(5,87.12)};
    \addplot[colorEfficientNet,mark=diamond]  coordinates{(1,83.20)(2,84.96)(3,85.20)(4,85.52)(5,85.60)};
    \addplot[colorViT,         mark=*]        coordinates{(1,83.12)(2,83.52)(3,83.92)(4,83.68)(5,83.68)};

    \Ann{south}{1}{80.9}{82.08}{2.79}{colorVGG}
    \Ann{north}{3}{84.5}{84.00}{2.56}{colorVGG}
    \Ann{south}{1}{81.8}{83.04}{3.92}{colorMobileNet}
    \Ann{north}{5}{87.1}{86.64}{3.51}{colorMobileNet}
    \Ann{north}{1}{84.6}{84.08}{2.50}{colorDenseNet}
    \Ann{south}{1}{80.6}{81.84}{2.62}{colorResNet}
    \Ann{north}{4}{88.0}{87.52}{2.57}{colorResNet}
    \Ann{north}{1}{83.7}{83.20}{2.46}{colorEfficientNet}
    \Ann{north}{1}{83.7}{83.12}{2.93}{colorViT}

    \SensLegend
  \end{axis}
  \end{tikzpicture}
  \end{adjustbox}
  \caption{Test Accuracy vs.\ $k$}
  \label{fig:tomato_acc}
\end{subfigure}\hfill
\begin{subfigure}[t]{0.49\textwidth}
  \begin{adjustbox}{max width=\linewidth}
  \begin{tikzpicture}
  \begin{axis}[sens plot,
    title={Tomato — F1-Score},
    ylabel={F1-Score (\%)},
    ymin=82.0, ymax=93.5]

    \addplot[colorVGG,         mark=x]        coordinates{(1,84.71)(2,86.04)(3,86.87)(4,86.96)(5,86.43)};
    \addplot[colorMobileNet,   mark=triangle] coordinates{(1,86.78)(2,87.84)(3,88.56)(4,88.93)(5,89.70)};
    \addplot[colorDenseNet,    mark=square]   coordinates{(1,87.52)(2,87.98)(3,87.91)(4,87.94)(5,87.95)};
    \addplot[colorResNet,      mark=triangle, mark options={rotate=180}]
                                              coordinates{(1,84.48)(2,88.08)(3,89.38)(4,90.01)(5,89.91)};
    \addplot[colorEfficientNet,mark=diamond]  coordinates{(1,86.19)(2,88.52)(3,88.58)(4,88.78)(5,88.99)};
    \addplot[colorViT,         mark=*]        coordinates{(1,86.71)(2,87.23)(3,87.37)(4,87.21)(5,87.26)};

    \Ann{south}{1}{83.5}{84.71}{2.67}{colorVGG}
    \Ann{north}{4}{87.4}{86.96}{1.91}{colorVGG}
    \Ann{south}{1}{85.5}{86.78}{3.99}{colorMobileNet}
    \Ann{north}{5}{90.2}{89.70}{2.80}{colorMobileNet}
    \Ann{north}{1}{88.0}{87.52}{1.94}{colorDenseNet}
    \Ann{south}{1}{83.2}{84.48}{1.92}{colorResNet}
    \Ann{north}{4}{90.5}{90.01}{1.84}{colorResNet}
    \Ann{south}{1}{84.9}{86.19}{2.21}{colorEfficientNet}
    \Ann{north}{5}{89.5}{88.99}{1.70}{colorEfficientNet}
    \Ann{north}{1}{87.2}{86.71}{2.77}{colorViT}

    \SensLegend
  \end{axis}
  \end{tikzpicture}
  \end{adjustbox}
  \caption{Test F1-Score vs.\ $k$}
  \label{fig:tomato_f1}
\end{subfigure}
\caption{Sensitivity analysis on \textbf{Tomato}.}
\label{fig:sensitivity_tomato}
\end{figure}

% ╔══════════════════════════════════════════════════════════════════════╗
% ║  FIGURE 5 — POTATO                                                  ║
% ╚══════════════════════════════════════════════════════════════════════╝
\begin{figure}[htbp]
\centering
\begin{subfigure}[t]{0.49\textwidth}
  \begin{adjustbox}{max width=\linewidth}
  \begin{tikzpicture}
  \begin{axis}[sens plot,
    title={Potato — Accuracy},
    ylabel={Accuracy (\%)},
    ymin=83.5, ymax=93.5]

    \addplot[colorVGG,         mark=x]        coordinates{(1,86.03)(2,87.18)(3,88.08)(4,87.61)(5,87.78)};
    \addplot[colorMobileNet,   mark=triangle] coordinates{(1,87.09)(2,88.03)(3,88.29)(4,88.50)(5,88.76)};
    \addplot[colorDenseNet,    mark=square]   coordinates{(1,86.79)(2,89.10)(3,89.66)(4,90.26)(5,90.00)};
    \addplot[colorResNet,      mark=triangle, mark options={rotate=180}]
                                              coordinates{(1,86.20)(2,88.42)(3,88.89)(4,88.76)(5,89.10)};
    \addplot[colorEfficientNet,mark=diamond]  coordinates{(1,85.94)(2,88.12)(3,88.46)(4,88.68)(5,88.97)};
    \addplot[colorViT,         mark=*]        coordinates{(1,90.30)(2,90.30)(3,90.43)(4,90.90)(5,90.60)};

    \Ann{south}{1}{84.8}{86.03}{1.67}{colorVGG}
    \Ann{north}{3}{88.6}{88.08}{1.03}{colorVGG}
    \Ann{south}{1}{85.8}{87.09}{1.67}{colorMobileNet}
    \Ann{north}{5}{89.2}{88.76}{0.72}{colorMobileNet}
    \Ann{south}{1}{85.6}{86.79}{1.88}{colorDenseNet}
    \Ann{north}{4}{90.7}{90.26}{2.01}{colorDenseNet}
    \Ann{south}{1}{85.0}{86.20}{1.02}{colorResNet}
    \Ann{north}{5}{89.6}{89.10}{0.79}{colorResNet}
    \Ann{south}{1}{84.7}{85.94}{0.76}{colorEfficientNet}
    \Ann{north}{5}{89.5}{88.97}{0.69}{colorEfficientNet}
    \Ann{north}{1}{90.8}{90.30}{1.51}{colorViT}
    \Ann{north}{4}{91.4}{90.90}{0.49}{colorViT}

    \SensLegend
  \end{axis}
  \end{tikzpicture}
  \end{adjustbox}
  \caption{Test Accuracy vs.\ $k$}
  \label{fig:potato_acc}
\end{subfigure}\hfill
\begin{subfigure}[t]{0.49\textwidth}
  \begin{adjustbox}{max width=\linewidth}
  \begin{tikzpicture}
  \begin{axis}[sens plot,
    title={Potato — F1-Score},
    ylabel={F1-Score (\%)},
    ymin=82.5, ymax=91.5]

    \addplot[colorVGG,         mark=x]        coordinates{(1,84.52)(2,86.07)(3,87.19)(4,86.68)(5,86.88)};
    \addplot[colorMobileNet,   mark=triangle] coordinates{(1,85.52)(2,86.89)(3,86.92)(4,87.21)(5,87.51)};
    \addplot[colorDenseNet,    mark=square]   coordinates{(1,84.90)(2,87.18)(3,87.64)(4,88.53)(5,88.38)};
    \addplot[colorResNet,      mark=triangle, mark options={rotate=180}]
                                              coordinates{(1,84.84)(2,87.53)(3,87.23)(4,87.50)(5,87.59)};
    \addplot[colorEfficientNet,mark=diamond]  coordinates{(1,84.67)(2,86.80)(3,87.13)(4,87.29)(5,87.57)};
    \addplot[colorViT,         mark=*]        coordinates{(1,88.06)(2,88.29)(3,88.53)(4,89.15)(5,88.94)};

    \Ann{south}{1}{83.3}{84.52}{2.20}{colorVGG}
    \Ann{north}{3}{87.7}{87.19}{1.23}{colorVGG}
    \Ann{south}{1}{84.3}{85.52}{1.94}{colorMobileNet}
    \Ann{north}{5}{88.0}{87.51}{0.81}{colorMobileNet}
    \Ann{south}{1}{83.7}{84.90}{1.90}{colorDenseNet}
    \Ann{north}{4}{89.0}{88.53}{2.22}{colorDenseNet}
    \Ann{south}{1}{83.6}{84.84}{0.68}{colorResNet}
    \Ann{south}{1}{83.4}{84.67}{0.94}{colorEfficientNet}
    \Ann{north}{5}{88.1}{87.57}{1.04}{colorEfficientNet}
    \Ann{north}{1}{88.6}{88.06}{1.31}{colorViT}
    \Ann{north}{4}{89.7}{89.15}{1.21}{colorViT}

    \SensLegend
  \end{axis}
  \end{tikzpicture}
  \end{adjustbox}
  \caption{Test F1-Score vs.\ $k$}
  \label{fig:potato_f1}
\end{subfigure}
\caption{Sensitivity analysis on \textbf{Potato}.}
\label{fig:sensitivity_potato}
\end{figure}

%  Requires in preamble: pgfplots, pgfplotsset{compat=1.18},
%    fillbetween, subcaption, xcolor
% ══════════════════════════════════════════════════════════════════════════

% ── Colors (already in your preamble — kept here for reference) ──────────
% \definecolor{colorViT}{HTML}{D62728}
% \definecolor{colorDenseNet}{HTML}{1F77B4}
% \definecolor{colorResNet}{HTML}{2CA02C}
% \definecolor{colorEfficientNet}{HTML}{FF7F0E}
% \definecolor{colorMobileNet}{HTML}{9467BD}
% \definecolor{colorVGG}{HTML}{7F7F7F}

% ── Shared axis style ────────────────────────────────────────────────────
\pgfplotsset{
  sens plot/.style={
    width=\linewidth,
    height=6.2cm,
    font=\small\rmfamily,
    xmin=0.55, xmax=5.45,
    xtick={1,2,3,4,5},
    xmajorgrids=true, ymajorgrids=true,
    grid style={dashed, gray!25},
    axis line style={gray!60},
    title style={font=\small\bfseries, yshift=-2pt},
    xlabel={Value of $k$\enspace{\footnotesize($k{=}1$: baseline)}},
    xlabel style={font=\small, yshift=3pt},
    ylabel style={font=\small},
    tick label style={font=\footnotesize},
    legend style={
      at={(0.5,-0.20)}, anchor=north,
      legend columns=3,
      font=\fontsize{6.8}{8}\selectfont,
      draw=gray!30, rounded corners=2pt,
      inner sep=3pt,
      /tikz/every even column/.append style={column sep=4pt},
    },
    line width=1.5pt,
    mark size=2.4pt,
    clip=false,   % allow annotations outside axis box
  }
}

% ── Annotation macro ─────────────────────────────────────────────────────
% \Ann{anchor}{x}{y}{val}{std}{color}
\newcommand{\Ann}[6]{%
  \node[anchor=#1, text=#6,
        font=\fontsize{5.8}{7}\selectfont\bfseries,
        fill=white, fill opacity=0.88, text opacity=1,
        inner sep=1pt, rounded corners=1.2pt]
    at (axis cs:#2,#3)
    {$#4$\\[-1.5pt]\normalfont$\pm\!#5$};}

% ── Shared legend entry (used in every subplot) ───────────────────────────
\newcommand{\SensLegend}{%
  \addlegendentry{VGG16}
  \addlegendentry{MobileNetV2}
  \addlegendentry{DenseNet121}
  \addlegendentry{ResNet101}
  \addlegendentry{EfficientNetB0}
  \addlegendentry{ViT-Base}
}


% ════════════════════════════════════════════════════════════════════════
%  HELPER: one subplot axis
%  Args: #1=ylabel text  #2=ymin  #3=ymax  #4=axis body (addplot + Ann)
% ════════════════════════════════════════════════════════════════════════
% (We don't use a separate macro for the axis body since pgfplots
%  needs the content inline — we just repeat the axis skeleton below.)


% ╔══════════════════════════════════════════════════════════════════════╗
% ║  FIGURE 1 — CIFAR-100                                               ║
% ╚══════════════════════════════════════════════════════════════════════╝
\begin{figure}[htbp]
\centering
\begin{subfigure}[t]{0.49\textwidth}
  \begin{adjustbox}{max width=\linewidth}
  \begin{tikzpicture}
  \begin{axis}[sens plot,
    title={CIFAR-100 — Accuracy},
    ylabel={Accuracy (\%)},
    ymin=69.5, ymax=91.5]

    \addplot[colorVGG,   mark=x]            coordinates{(1,71.57)(2,72.10)(3,72.34)(4,72.48)(5,72.46)};
    \addplot[colorMobileNet, mark=triangle] coordinates{(1,73.23)(2,74.24)(3,74.68)(4,74.92)(5,75.06)};
    \addplot[colorDenseNet,  mark=square]   coordinates{(1,75.39)(2,78.10)(3,78.73)(4,78.93)(5,79.23)};
    \addplot[colorResNet,    mark=triangle, mark options={rotate=180}]
                                            coordinates{(1,75.84)(2,78.08)(3,78.60)(4,78.82)(5,78.80)};
    \addplot[colorEfficientNet, mark=diamond]
                                            coordinates{(1,76.76)(2,79.85)(3,80.51)(4,80.92)(5,81.15)};
    \addplot[colorViT,   mark=*]            coordinates{(1,87.19)(2,88.68)(3,89.12)(4,89.39)(5,89.64)};

    % NOTE: a partial "std band" overlay for VGG16 was removed here. It was the
    % only band in the whole file and \closedcycle fills down to the axis rather
    % than between the two curves, so it rendered as a shaded wedge under VGG16
    % alone. The per-point std is already given by the \Ann annotations.

    % ── Annotations: baseline k=1 & peak ──
    \Ann{south}{1}{70.9}{71.57}{0.41}{colorVGG}
    \Ann{north}{1}{73.6}{73.23}{0.41}{colorMobileNet}
    \Ann{north}{5}{75.4}{75.06}{0.15}{colorMobileNet}
    \Ann{south}{1}{74.7}{75.39}{0.20}{colorDenseNet}
    \Ann{north}{5}{79.6}{79.23}{0.30}{colorDenseNet}
    \Ann{north}{1}{76.4}{75.84}{1.54}{colorResNet}
    \Ann{north}{1}{77.2}{76.76}{0.51}{colorEfficientNet}
    \Ann{north}{5}{81.6}{81.15}{0.31}{colorEfficientNet}
    \Ann{north}{1}{87.7}{87.19}{0.65}{colorViT}
    \Ann{north}{5}{90.1}{89.64}{0.28}{colorViT}

    \SensLegend
  \end{axis}
  \end{tikzpicture}
  \end{adjustbox}
  \caption{Test Accuracy vs.\ $k$}
  \label{fig:cifar100_acc}
\end{subfigure}\hfill
\begin{subfigure}[t]{0.49\textwidth}
  \begin{adjustbox}{max width=\linewidth}
  \begin{tikzpicture}
  \begin{axis}[sens plot,
    title={CIFAR-100 — F1-Score},
    ylabel={F1-Score (\%)},
    ymin=69.5, ymax=91.5]

    \addplot[colorVGG,         mark=x]            coordinates{(1,71.59)(2,72.11)(3,72.36)(4,72.50)(5,72.48)};
    \addplot[colorMobileNet,   mark=triangle]     coordinates{(1,73.25)(2,74.27)(3,74.71)(4,74.95)(5,75.04)};
    \addplot[colorDenseNet,    mark=square]       coordinates{(1,75.41)(2,78.09)(3,78.72)(4,78.92)(5,79.21)};
    \addplot[colorResNet,      mark=triangle, mark options={rotate=180}]
                                                  coordinates{(1,75.92)(2,78.15)(3,78.67)(4,78.87)(5,78.87)};
    \addplot[colorEfficientNet, mark=diamond]     coordinates{(1,76.75)(2,79.82)(3,80.46)(4,80.88)(5,81.10)};
    \addplot[colorViT,         mark=*]            coordinates{(1,87.01)(2,88.52)(3,88.95)(4,89.25)(5,89.52)};

    \Ann{south}{1}{70.9}{71.59}{0.42}{colorVGG}
    \Ann{north}{1}{73.7}{73.25}{0.50}{colorMobileNet}
    \Ann{north}{5}{75.4}{75.04}{0.18}{colorMobileNet}
    \Ann{south}{1}{74.7}{75.41}{0.28}{colorDenseNet}
    \Ann{north}{5}{79.6}{79.21}{0.28}{colorDenseNet}
    \Ann{north}{1}{76.5}{75.92}{1.52}{colorResNet}
    \Ann{north}{1}{77.2}{76.75}{0.52}{colorEfficientNet}
    \Ann{north}{5}{81.6}{81.10}{0.34}{colorEfficientNet}
    \Ann{north}{1}{87.6}{87.01}{0.69}{colorViT}
    \Ann{north}{5}{90.0}{89.52}{0.30}{colorViT}

    \SensLegend
  \end{axis}
  \end{tikzpicture}
  \end{adjustbox}
  \caption{Test F1-Score vs.\ $k$}
  \label{fig:cifar100_f1}
\end{subfigure}
\caption{Sensitivity analysis on \textbf{CIFAR-100}: effect of checkpoint-averaging
  parameter $k$ on Test Accuracy and F1-Score. $k{=}1$ is the single-checkpoint
  baseline. Annotated values show mean\,$\pm$\,std over five independent runs.}
\label{fig:sensitivity_cifar100}
\end{figure}

% ╔══════════════════════════════════════════════════════════════════════╗
% ║  FIGURE 2 — TINY IMAGENET                                           ║
% ╚══════════════════════════════════════════════════════════════════════╝
\begin{figure}[htbp]
\centering
\begin{subfigure}[t]{0.49\textwidth}
  \begin{adjustbox}{max width=\linewidth}
  \begin{tikzpicture}
  \begin{axis}[sens plot,
    title={Tiny ImageNet — Validation Accuracy},
    ylabel={Accuracy (\%)},
    ymin=65.5, ymax=93.0]

    \addplot[colorVGG,         mark=x]        coordinates{(1,67.54)(2,68.04)(3,68.12)(4,68.35)(5,68.34)};
    \addplot[colorMobileNet,   mark=triangle] coordinates{(1,72.90)(2,73.39)(3,73.40)(4,73.46)(5,73.53)};
    \addplot[colorDenseNet,    mark=square]   coordinates{(1,76.82)(2,77.74)(3,77.88)(4,78.06)(5,78.04)};
    \addplot[colorResNet,      mark=triangle, mark options={rotate=180}]
                                              coordinates{(1,77.50)(2,80.41)(3,80.86)(4,81.23)(5,81.34)};
    \addplot[colorEfficientNet,mark=diamond]  coordinates{(1,80.32)(2,80.64)(3,80.72)(4,80.81)(5,80.78)};
    \addplot[colorViT,         mark=*]        coordinates{(1,89.68)(2,90.96)(3,91.31)(4,91.43)(5,91.53)};

    \Ann{south}{1}{66.5}{67.54}{0.85}{colorVGG}
    \Ann{south}{1}{71.7}{72.90}{0.54}{colorMobileNet}
    \Ann{north}{5}{73.9}{73.53}{0.21}{colorMobileNet}
    \Ann{south}{1}{75.8}{76.82}{0.61}{colorDenseNet}
    \Ann{north}{5}{78.5}{78.04}{0.22}{colorDenseNet}
    \Ann{south}{1}{76.4}{77.50}{0.60}{colorResNet}
    \Ann{north}{5}{81.8}{81.34}{0.32}{colorResNet}
    \Ann{north}{1}{80.9}{80.32}{0.32}{colorEfficientNet}
    \Ann{north}{1}{90.2}{89.68}{0.34}{colorViT}
    \Ann{north}{5}{92.0}{91.53}{0.12}{colorViT}

    \SensLegend
  \end{axis}
  \end{tikzpicture}
  \end{adjustbox}
  \caption{Validation Accuracy vs.\ $k$}
  \label{fig:tinyimagenet_acc}
\end{subfigure}\hfill
\begin{subfigure}[t]{0.49\textwidth}
  \begin{adjustbox}{max width=\linewidth}
  \begin{tikzpicture}
  \begin{axis}[sens plot,
    title={Tiny ImageNet — Validation Macro-F1},
    ylabel={F1-Score (\%)},
    ymin=65.5, ymax=93.0]

    \addplot[colorVGG,         mark=x]        coordinates{(1,67.36)(2,67.85)(3,67.91)(4,68.14)(5,68.13)};
    \addplot[colorMobileNet,   mark=triangle] coordinates{(1,72.74)(2,73.24)(3,73.25)(4,73.31)(5,73.38)};
    \addplot[colorDenseNet,    mark=square]   coordinates{(1,76.81)(2,77.73)(3,77.85)(4,78.03)(5,78.01)};
    \addplot[colorResNet,      mark=triangle, mark options={rotate=180}]
                                              coordinates{(1,77.58)(2,80.41)(3,80.85)(4,81.23)(5,81.34)};
    \addplot[colorEfficientNet,mark=diamond]  coordinates{(1,80.25)(2,80.58)(3,80.66)(4,80.75)(5,80.72)};
    \addplot[colorViT,         mark=*]        coordinates{(1,89.69)(2,90.96)(3,91.31)(4,91.43)(5,91.52)};

    \Ann{south}{1}{66.3}{67.36}{0.87}{colorVGG}
    \Ann{south}{1}{71.6}{72.74}{0.54}{colorMobileNet}
    \Ann{north}{5}{73.8}{73.38}{0.21}{colorMobileNet}
    \Ann{south}{1}{75.7}{76.81}{0.59}{colorDenseNet}
    \Ann{north}{5}{78.4}{78.01}{0.21}{colorDenseNet}
    \Ann{south}{1}{76.4}{77.58}{0.51}{colorResNet}
    \Ann{north}{5}{81.8}{81.34}{0.34}{colorResNet}
    \Ann{north}{1}{80.8}{80.25}{0.32}{colorEfficientNet}
    \Ann{north}{1}{90.2}{89.69}{0.33}{colorViT}
    \Ann{north}{5}{92.0}{91.52}{0.12}{colorViT}

    \SensLegend
  \end{axis}
  \end{tikzpicture}
  \end{adjustbox}
  \caption{Validation Macro-F1 vs.\ $k$}
  \label{fig:tinyimagenet_f1}
\end{subfigure}
\caption{Sensitivity analysis on \textbf{Tiny ImageNet}. The official validation split is used as the held-out evaluation set because test labels are unavailable, so both panels report validation metrics.}
\label{fig:sensitivity_tinyimagenet}
\end{figure}

% ╔══════════════════════════════════════════════════════════════════════╗
% ║  FIGURE 3 — BURMESE GRAPE                                           ║
% ╚══════════════════════════════════════════════════════════════════════╝
\begin{figure}[htbp]
\centering
\begin{subfigure}[t]{0.49\textwidth}
  \begin{adjustbox}{max width=\linewidth}
  \begin{tikzpicture}
  \begin{axis}[sens plot,
    title={Burmese Grape — Accuracy},
    ylabel={Accuracy (\%)},
    ymin=79.5, ymax=90.5]

    \addplot[colorVGG,         mark=x]        coordinates{(1,83.76)(2,85.30)(3,85.43)(4,85.77)(5,85.94)};
    \addplot[colorMobileNet,   mark=triangle] coordinates{(1,83.89)(2,85.94)(3,86.67)(4,87.19)(5,87.01)};
    \addplot[colorDenseNet,    mark=square]   coordinates{(1,84.62)(2,84.96)(3,85.00)(4,84.96)(5,85.17)};
    \addplot[colorResNet,      mark=triangle, mark options={rotate=180}]
                                              coordinates{(1,85.38)(2,85.13)(3,85.17)(4,85.68)(5,85.60)};
    \addplot[colorEfficientNet,mark=diamond]  coordinates{(1,82.74)(2,85.56)(3,85.77)(4,86.54)(5,86.45)};
    \addplot[colorViT,         mark=*]        coordinates{(1,86.98)(2,86.58)(3,87.05)(4,86.79)(5,87.01)};

    \Ann{south}{1}{82.6}{83.76}{2.04}{colorVGG}
    \Ann{north}{5}{86.3}{85.94}{0.98}{colorVGG}
    \Ann{south}{1}{82.7}{83.89}{1.54}{colorMobileNet}
    \Ann{north}{4}{87.6}{87.19}{1.54}{colorMobileNet}
    \Ann{north}{1}{85.0}{84.62}{1.62}{colorDenseNet}
    \Ann{south}{1}{84.3}{85.38}{0.94}{colorResNet}
    \Ann{south}{1}{81.5}{82.74}{2.09}{colorEfficientNet}
    \Ann{north}{4}{87.0}{86.54}{1.05}{colorEfficientNet}
    \Ann{north}{1}{87.5}{86.98}{1.92}{colorViT}

    \SensLegend
  \end{axis}
  \end{tikzpicture}
  \end{adjustbox}
  \caption{Test Accuracy vs.\ $k$}
  \label{fig:grape_acc}
\end{subfigure}\hfill
\begin{subfigure}[t]{0.49\textwidth}
  \begin{adjustbox}{max width=\linewidth}
  \begin{tikzpicture}
  \begin{axis}[sens plot,
    title={Burmese Grape — F1-Score},
    ylabel={F1-Score (\%)},
    ymin=78.0, ymax=90.0]

    \addplot[colorVGG,         mark=x]        coordinates{(1,81.83)(2,83.97)(3,83.69)(4,84.08)(5,84.41)};
    \addplot[colorMobileNet,   mark=triangle] coordinates{(1,81.85)(2,84.55)(3,85.42)(4,85.64)(5,86.08)};
    \addplot[colorDenseNet,    mark=square]   coordinates{(1,82.44)(2,83.01)(3,83.01)(4,82.84)(5,83.06)};
    \addplot[colorResNet,      mark=triangle, mark options={rotate=180}]
                                              coordinates{(1,84.21)(2,84.02)(3,84.08)(4,84.54)(5,84.40)};
    \addplot[colorEfficientNet,mark=diamond]  coordinates{(1,81.17)(2,84.39)(3,84.84)(4,85.49)(5,85.51)};
    \addplot[colorViT,         mark=*]        coordinates{(1,84.85)(2,85.70)(3,86.00)(4,85.66)(5,85.89)};

    \Ann{south}{1}{80.6}{81.83}{3.42}{colorVGG}
    \Ann{north}{5}{84.8}{84.41}{1.42}{colorVGG}
    \Ann{south}{1}{80.6}{81.85}{1.62}{colorMobileNet}
    \Ann{north}{5}{86.5}{86.08}{1.39}{colorMobileNet}
    \Ann{north}{1}{82.9}{82.44}{1.46}{colorDenseNet}
    \Ann{north}{1}{84.7}{84.21}{1.17}{colorResNet}
    \Ann{south}{1}{79.9}{81.17}{2.70}{colorEfficientNet}
    \Ann{north}{5}{85.9}{85.51}{0.90}{colorEfficientNet}
    \Ann{north}{1}{85.3}{84.85}{1.49}{colorViT}
    \Ann{north}{3}{86.5}{86.00}{1.87}{colorViT}

    \SensLegend
  \end{axis}
  \end{tikzpicture}
  \end{adjustbox}
  \caption{Test F1-Score vs.\ $k$}
  \label{fig:grape_f1}
\end{subfigure}
\caption{Sensitivity analysis on \textbf{Burmese Grape}.}
\label{fig:sensitivity_grape}
\end{figure}

% ╔══════════════════════════════════════════════════════════════════════╗
% ║  FIGURE 4 — TOMATO                                                  ║
% ╚══════════════════════════════════════════════════════════════════════╝
\begin{figure}[htbp]
\centering
\begin{subfigure}[t]{0.49\textwidth}
  \begin{adjustbox}{max width=\linewidth}
  \begin{tikzpicture}
  \begin{axis}[sens plot,
    title={Tomato — Accuracy},
    ylabel={Accuracy (\%)},
    ymin=78.5, ymax=91.0]

    \addplot[colorVGG,         mark=x]        coordinates{(1,82.08)(2,83.36)(3,84.00)(4,84.00)(5,83.20)};
    \addplot[colorMobileNet,   mark=triangle] coordinates{(1,83.04)(2,84.40)(3,85.12)(4,85.92)(5,86.64)};
    \addplot[colorDenseNet,    mark=square]   coordinates{(1,84.08)(2,84.56)(3,84.48)(4,84.48)(5,84.40)};
    \addplot[colorResNet,      mark=triangle, mark options={rotate=180}]
                                              coordinates{(1,81.84)(2,85.36)(3,86.64)(4,87.52)(5,87.12)};
    \addplot[colorEfficientNet,mark=diamond]  coordinates{(1,83.20)(2,84.96)(3,85.20)(4,85.52)(5,85.60)};
    \addplot[colorViT,         mark=*]        coordinates{(1,83.12)(2,83.52)(3,83.92)(4,83.68)(5,83.68)};

    \Ann{south}{1}{80.9}{82.08}{2.79}{colorVGG}
    \Ann{north}{3}{84.5}{84.00}{2.56}{colorVGG}
    \Ann{south}{1}{81.8}{83.04}{3.92}{colorMobileNet}
    \Ann{north}{5}{87.1}{86.64}{3.51}{colorMobileNet}
    \Ann{north}{1}{84.6}{84.08}{2.50}{colorDenseNet}
    \Ann{south}{1}{80.6}{81.84}{2.62}{colorResNet}
    \Ann{north}{4}{88.0}{87.52}{2.57}{colorResNet}
    \Ann{north}{1}{83.7}{83.20}{2.46}{colorEfficientNet}
    \Ann{north}{1}{83.7}{83.12}{2.93}{colorViT}

    \SensLegend
  \end{axis}
  \end{tikzpicture}
  \end{adjustbox}
  \caption{Test Accuracy vs.\ $k$}
  \label{fig:tomato_acc}
\end{subfigure}\hfill
\begin{subfigure}[t]{0.49\textwidth}
  \begin{adjustbox}{max width=\linewidth}
  \begin{tikzpicture}
  \begin{axis}[sens plot,
    title={Tomato — F1-Score},
    ylabel={F1-Score (\%)},
    ymin=82.0, ymax=93.5]

    \addplot[colorVGG,         mark=x]        coordinates{(1,84.71)(2,86.04)(3,86.87)(4,86.96)(5,86.43)};
    \addplot[colorMobileNet,   mark=triangle] coordinates{(1,86.78)(2,87.84)(3,88.56)(4,88.93)(5,89.70)};
    \addplot[colorDenseNet,    mark=square]   coordinates{(1,87.52)(2,87.98)(3,87.91)(4,87.94)(5,87.95)};
    \addplot[colorResNet,      mark=triangle, mark options={rotate=180}]
                                              coordinates{(1,84.48)(2,88.08)(3,89.38)(4,90.01)(5,89.91)};
    \addplot[colorEfficientNet,mark=diamond]  coordinates{(1,86.19)(2,88.52)(3,88.58)(4,88.78)(5,88.99)};
    \addplot[colorViT,         mark=*]        coordinates{(1,86.71)(2,87.23)(3,87.37)(4,87.21)(5,87.26)};

    \Ann{south}{1}{83.5}{84.71}{2.67}{colorVGG}
    \Ann{north}{4}{87.4}{86.96}{1.91}{colorVGG}
    \Ann{south}{1}{85.5}{86.78}{3.99}{colorMobileNet}
    \Ann{north}{5}{90.2}{89.70}{2.80}{colorMobileNet}
    \Ann{north}{1}{88.0}{87.52}{1.94}{colorDenseNet}
    \Ann{south}{1}{83.2}{84.48}{1.92}{colorResNet}
    \Ann{north}{4}{90.5}{90.01}{1.84}{colorResNet}
    \Ann{south}{1}{84.9}{86.19}{2.21}{colorEfficientNet}
    \Ann{north}{5}{89.5}{88.99}{1.70}{colorEfficientNet}
    \Ann{north}{1}{87.2}{86.71}{2.77}{colorViT}

    \SensLegend
  \end{axis}
  \end{tikzpicture}
  \end{adjustbox}
  \caption{Test F1-Score vs.\ $k$}
  \label{fig:tomato_f1}
\end{subfigure}
\caption{Sensitivity analysis on \textbf{Tomato}.}
\label{fig:sensitivity_tomato}
\end{figure}

% ╔══════════════════════════════════════════════════════════════════════╗
% ║  FIGURE 5 — POTATO                                                  ║
% ╚══════════════════════════════════════════════════════════════════════╝
\begin{figure}[htbp]
\centering
\begin{subfigure}[t]{0.49\textwidth}
  \begin{adjustbox}{max width=\linewidth}
  \begin{tikzpicture}
  \begin{axis}[sens plot,
    title={Potato — Accuracy},
    ylabel={Accuracy (\%)},
    ymin=83.5, ymax=93.5]

    \addplot[colorVGG,         mark=x]        coordinates{(1,86.03)(2,87.18)(3,88.08)(4,87.61)(5,87.78)};
    \addplot[colorMobileNet,   mark=triangle] coordinates{(1,87.09)(2,88.03)(3,88.29)(4,88.50)(5,88.76)};
    \addplot[colorDenseNet,    mark=square]   coordinates{(1,86.79)(2,89.10)(3,89.66)(4,90.26)(5,90.00)};
    \addplot[colorResNet,      mark=triangle, mark options={rotate=180}]
                                              coordinates{(1,86.20)(2,88.42)(3,88.89)(4,88.76)(5,89.10)};
    \addplot[colorEfficientNet,mark=diamond]  coordinates{(1,85.94)(2,88.12)(3,88.46)(4,88.68)(5,88.97)};
    \addplot[colorViT,         mark=*]        coordinates{(1,90.30)(2,90.30)(3,90.43)(4,90.90)(5,90.60)};

    \Ann{south}{1}{84.8}{86.03}{1.67}{colorVGG}
    \Ann{north}{3}{88.6}{88.08}{1.03}{colorVGG}
    \Ann{south}{1}{85.8}{87.09}{1.67}{colorMobileNet}
    \Ann{north}{5}{89.2}{88.76}{0.72}{colorMobileNet}
    \Ann{south}{1}{85.6}{86.79}{1.88}{colorDenseNet}
    \Ann{north}{4}{90.7}{90.26}{2.01}{colorDenseNet}
    \Ann{south}{1}{85.0}{86.20}{1.02}{colorResNet}
    \Ann{north}{5}{89.6}{89.10}{0.79}{colorResNet}
    \Ann{south}{1}{84.7}{85.94}{0.76}{colorEfficientNet}
    \Ann{north}{5}{89.5}{88.97}{0.69}{colorEfficientNet}
    \Ann{north}{1}{90.8}{90.30}{1.51}{colorViT}
    \Ann{north}{4}{91.4}{90.90}{0.49}{colorViT}

    \SensLegend
  \end{axis}
  \end{tikzpicture}
  \end{adjustbox}
  \caption{Test Accuracy vs.\ $k$}
  \label{fig:potato_acc}
\end{subfigure}\hfill
\begin{subfigure}[t]{0.49\textwidth}
  \begin{adjustbox}{max width=\linewidth}
  \begin{tikzpicture}
  \begin{axis}[sens plot,
    title={Potato — F1-Score},
    ylabel={F1-Score (\%)},
    ymin=82.5, ymax=91.5]

    \addplot[colorVGG,         mark=x]        coordinates{(1,84.52)(2,86.07)(3,87.19)(4,86.68)(5,86.88)};
    \addplot[colorMobileNet,   mark=triangle] coordinates{(1,85.52)(2,86.89)(3,86.92)(4,87.21)(5,87.51)};
    \addplot[colorDenseNet,    mark=square]   coordinates{(1,84.90)(2,87.18)(3,87.64)(4,88.53)(5,88.38)};
    \addplot[colorResNet,      mark=triangle, mark options={rotate=180}]
                                              coordinates{(1,84.84)(2,87.53)(3,87.23)(4,87.50)(5,87.59)};
    \addplot[colorEfficientNet,mark=diamond]  coordinates{(1,84.67)(2,86.80)(3,87.13)(4,87.29)(5,87.57)};
    \addplot[colorViT,         mark=*]        coordinates{(1,88.06)(2,88.29)(3,88.53)(4,89.15)(5,88.94)};

    \Ann{south}{1}{83.3}{84.52}{2.20}{colorVGG}
    \Ann{north}{3}{87.7}{87.19}{1.23}{colorVGG}
    \Ann{south}{1}{84.3}{85.52}{1.94}{colorMobileNet}
    \Ann{north}{5}{88.0}{87.51}{0.81}{colorMobileNet}
    \Ann{south}{1}{83.7}{84.90}{1.90}{colorDenseNet}
    \Ann{north}{4}{89.0}{88.53}{2.22}{colorDenseNet}
    \Ann{south}{1}{83.6}{84.84}{0.68}{colorResNet}
    \Ann{south}{1}{83.4}{84.67}{0.94}{colorEfficientNet}
    \Ann{north}{5}{88.1}{87.57}{1.04}{colorEfficientNet}
    \Ann{north}{1}{88.6}{88.06}{1.31}{colorViT}
    \Ann{north}{4}{89.7}{89.15}{1.21}{colorViT}

    \SensLegend
  \end{axis}
  \end{tikzpicture}
  \end{adjustbox}
  \caption{Test F1-Score vs.\ $k$}
  \label{fig:potato_f1}
\end{subfigure}
\caption{Sensitivity analysis on \textbf{Potato}.}
\label{fig:sensitivity_potato}
\end{figure}

%  Requires in preamble: pgfplots, pgfplotsset{compat=1.18},
%    fillbetween, subcaption, xcolor
% ══════════════════════════════════════════════════════════════════════════

% ── Colors (already in your preamble — kept here for reference) ──────────
% \definecolor{colorViT}{HTML}{D62728}
% \definecolor{colorDenseNet}{HTML}{1F77B4}
% \definecolor{colorResNet}{HTML}{2CA02C}
% \definecolor{colorEfficientNet}{HTML}{FF7F0E}
% \definecolor{colorMobileNet}{HTML}{9467BD}
% \definecolor{colorVGG}{HTML}{7F7F7F}

% ── Shared axis style ────────────────────────────────────────────────────
\pgfplotsset{
  sens plot/.style={
    width=\linewidth,
    height=6.2cm,
    font=\small\rmfamily,
    xmin=0.55, xmax=5.45,
    xtick={1,2,3,4,5},
    xmajorgrids=true, ymajorgrids=true,
    grid style={dashed, gray!25},
    axis line style={gray!60},
    title style={font=\small\bfseries, yshift=-2pt},
    xlabel={Value of $k$\enspace{\footnotesize($k{=}1$: baseline)}},
    xlabel style={font=\small, yshift=3pt},
    ylabel style={font=\small},
    tick label style={font=\footnotesize},
    legend style={
      at={(0.5,-0.20)}, anchor=north,
      legend columns=3,
      font=\fontsize{6.8}{8}\selectfont,
      draw=gray!30, rounded corners=2pt,
      inner sep=3pt,
      /tikz/every even column/.append style={column sep=4pt},
    },
    line width=1.5pt,
    mark size=2.4pt,
    clip=false,   % allow annotations outside axis box
  }
}

% ── Annotation macro ─────────────────────────────────────────────────────
% \Ann{anchor}{x}{y}{val}{std}{color}
\newcommand{\Ann}[6]{%
  \node[anchor=#1, text=#6,
        font=\fontsize{5.8}{7}\selectfont\bfseries,
        fill=white, fill opacity=0.88, text opacity=1,
        inner sep=1pt, rounded corners=1.2pt]
    at (axis cs:#2,#3)
    {$#4$\\[-1.5pt]\normalfont$\pm\!#5$};}

% ── Shared legend entry (used in every subplot) ───────────────────────────
\newcommand{\SensLegend}{%
  \addlegendentry{VGG16}
  \addlegendentry{MobileNetV2}
  \addlegendentry{DenseNet121}
  \addlegendentry{ResNet101}
  \addlegendentry{EfficientNetB0}
  \addlegendentry{ViT-Base}
}


% ════════════════════════════════════════════════════════════════════════
%  HELPER: one subplot axis
%  Args: #1=ylabel text  #2=ymin  #3=ymax  #4=axis body (addplot + Ann)
% ════════════════════════════════════════════════════════════════════════
% (We don't use a separate macro for the axis body since pgfplots
%  needs the content inline — we just repeat the axis skeleton below.)


% ╔══════════════════════════════════════════════════════════════════════╗
% ║  FIGURE 1 — CIFAR-100                                               ║
% ╚══════════════════════════════════════════════════════════════════════╝
\begin{figure}[htbp]
\centering
\begin{subfigure}[t]{0.49\textwidth}
  \begin{adjustbox}{max width=\linewidth}
  \begin{tikzpicture}
  \begin{axis}[sens plot,
    title={CIFAR-100 — Accuracy},
    ylabel={Accuracy (\%)},
    ymin=69.5, ymax=91.5]

    \addplot[colorVGG,   mark=x]            coordinates{(1,71.57)(2,72.10)(3,72.34)(4,72.48)(5,72.46)};
    \addplot[colorMobileNet, mark=triangle] coordinates{(1,73.23)(2,74.24)(3,74.68)(4,74.92)(5,75.06)};
    \addplot[colorDenseNet,  mark=square]   coordinates{(1,75.39)(2,78.10)(3,78.73)(4,78.93)(5,79.23)};
    \addplot[colorResNet,    mark=triangle, mark options={rotate=180}]
                                            coordinates{(1,75.84)(2,78.08)(3,78.60)(4,78.82)(5,78.80)};
    \addplot[colorEfficientNet, mark=diamond]
                                            coordinates{(1,76.76)(2,79.85)(3,80.51)(4,80.92)(5,81.15)};
    \addplot[colorViT,   mark=*]            coordinates{(1,87.19)(2,88.68)(3,89.12)(4,89.39)(5,89.64)};

    % NOTE: a partial "std band" overlay for VGG16 was removed here. It was the
    % only band in the whole file and \closedcycle fills down to the axis rather
    % than between the two curves, so it rendered as a shaded wedge under VGG16
    % alone. The per-point std is already given by the \Ann annotations.

    % ── Annotations: baseline k=1 & peak ──
    \Ann{south}{1}{70.9}{71.57}{0.41}{colorVGG}
    \Ann{north}{1}{73.6}{73.23}{0.41}{colorMobileNet}
    \Ann{north}{5}{75.4}{75.06}{0.15}{colorMobileNet}
    \Ann{south}{1}{74.7}{75.39}{0.20}{colorDenseNet}
    \Ann{north}{5}{79.6}{79.23}{0.30}{colorDenseNet}
    \Ann{north}{1}{76.4}{75.84}{1.54}{colorResNet}
    \Ann{north}{1}{77.2}{76.76}{0.51}{colorEfficientNet}
    \Ann{north}{5}{81.6}{81.15}{0.31}{colorEfficientNet}
    \Ann{north}{1}{87.7}{87.19}{0.65}{colorViT}
    \Ann{north}{5}{90.1}{89.64}{0.28}{colorViT}

    \SensLegend
  \end{axis}
  \end{tikzpicture}
  \end{adjustbox}
  \caption{Test Accuracy vs.\ $k$}
  \label{fig:cifar100_acc}
\end{subfigure}\hfill
\begin{subfigure}[t]{0.49\textwidth}
  \begin{adjustbox}{max width=\linewidth}
  \begin{tikzpicture}
  \begin{axis}[sens plot,
    title={CIFAR-100 — F1-Score},
    ylabel={F1-Score (\%)},
    ymin=69.5, ymax=91.5]

    \addplot[colorVGG,         mark=x]            coordinates{(1,71.59)(2,72.11)(3,72.36)(4,72.50)(5,72.48)};
    \addplot[colorMobileNet,   mark=triangle]     coordinates{(1,73.25)(2,74.27)(3,74.71)(4,74.95)(5,75.04)};
    \addplot[colorDenseNet,    mark=square]       coordinates{(1,75.41)(2,78.09)(3,78.72)(4,78.92)(5,79.21)};
    \addplot[colorResNet,      mark=triangle, mark options={rotate=180}]
                                                  coordinates{(1,75.92)(2,78.15)(3,78.67)(4,78.87)(5,78.87)};
    \addplot[colorEfficientNet, mark=diamond]     coordinates{(1,76.75)(2,79.82)(3,80.46)(4,80.88)(5,81.10)};
    \addplot[colorViT,         mark=*]            coordinates{(1,87.01)(2,88.52)(3,88.95)(4,89.25)(5,89.52)};

    \Ann{south}{1}{70.9}{71.59}{0.42}{colorVGG}
    \Ann{north}{1}{73.7}{73.25}{0.50}{colorMobileNet}
    \Ann{north}{5}{75.4}{75.04}{0.18}{colorMobileNet}
    \Ann{south}{1}{74.7}{75.41}{0.28}{colorDenseNet}
    \Ann{north}{5}{79.6}{79.21}{0.28}{colorDenseNet}
    \Ann{north}{1}{76.5}{75.92}{1.52}{colorResNet}
    \Ann{north}{1}{77.2}{76.75}{0.52}{colorEfficientNet}
    \Ann{north}{5}{81.6}{81.10}{0.34}{colorEfficientNet}
    \Ann{north}{1}{87.6}{87.01}{0.69}{colorViT}
    \Ann{north}{5}{90.0}{89.52}{0.30}{colorViT}

    \SensLegend
  \end{axis}
  \end{tikzpicture}
  \end{adjustbox}
  \caption{Test F1-Score vs.\ $k$}
  \label{fig:cifar100_f1}
\end{subfigure}
\caption{Sensitivity analysis on \textbf{CIFAR-100}: effect of checkpoint-averaging
  parameter $k$ on Test Accuracy and F1-Score. $k{=}1$ is the single-checkpoint
  baseline. Annotated values show mean\,$\pm$\,std over five independent runs.}
\label{fig:sensitivity_cifar100}
\end{figure}

% ╔══════════════════════════════════════════════════════════════════════╗
% ║  FIGURE 2 — TINY IMAGENET                                           ║
% ╚══════════════════════════════════════════════════════════════════════╝
\begin{figure}[htbp]
\centering
\begin{subfigure}[t]{0.49\textwidth}
  \begin{adjustbox}{max width=\linewidth}
  \begin{tikzpicture}
  \begin{axis}[sens plot,
    title={Tiny ImageNet — Validation Accuracy},
    ylabel={Accuracy (\%)},
    ymin=65.5, ymax=93.0]

    \addplot[colorVGG,         mark=x]        coordinates{(1,67.54)(2,68.04)(3,68.12)(4,68.35)(5,68.34)};
    \addplot[colorMobileNet,   mark=triangle] coordinates{(1,72.90)(2,73.39)(3,73.40)(4,73.46)(5,73.53)};
    \addplot[colorDenseNet,    mark=square]   coordinates{(1,76.82)(2,77.74)(3,77.88)(4,78.06)(5,78.04)};
    \addplot[colorResNet,      mark=triangle, mark options={rotate=180}]
                                              coordinates{(1,77.50)(2,80.41)(3,80.86)(4,81.23)(5,81.34)};
    \addplot[colorEfficientNet,mark=diamond]  coordinates{(1,80.32)(2,80.64)(3,80.72)(4,80.81)(5,80.78)};
    \addplot[colorViT,         mark=*]        coordinates{(1,89.68)(2,90.96)(3,91.31)(4,91.43)(5,91.53)};

    \Ann{south}{1}{66.5}{67.54}{0.85}{colorVGG}
    \Ann{south}{1}{71.7}{72.90}{0.54}{colorMobileNet}
    \Ann{north}{5}{73.9}{73.53}{0.21}{colorMobileNet}
    \Ann{south}{1}{75.8}{76.82}{0.61}{colorDenseNet}
    \Ann{north}{5}{78.5}{78.04}{0.22}{colorDenseNet}
    \Ann{south}{1}{76.4}{77.50}{0.60}{colorResNet}
    \Ann{north}{5}{81.8}{81.34}{0.32}{colorResNet}
    \Ann{north}{1}{80.9}{80.32}{0.32}{colorEfficientNet}
    \Ann{north}{1}{90.2}{89.68}{0.34}{colorViT}
    \Ann{north}{5}{92.0}{91.53}{0.12}{colorViT}

    \SensLegend
  \end{axis}
  \end{tikzpicture}
  \end{adjustbox}
  \caption{Validation Accuracy vs.\ $k$}
  \label{fig:tinyimagenet_acc}
\end{subfigure}\hfill
\begin{subfigure}[t]{0.49\textwidth}
  \begin{adjustbox}{max width=\linewidth}
  \begin{tikzpicture}
  \begin{axis}[sens plot,
    title={Tiny ImageNet — Validation Macro-F1},
    ylabel={F1-Score (\%)},
    ymin=65.5, ymax=93.0]

    \addplot[colorVGG,         mark=x]        coordinates{(1,67.36)(2,67.85)(3,67.91)(4,68.14)(5,68.13)};
    \addplot[colorMobileNet,   mark=triangle] coordinates{(1,72.74)(2,73.24)(3,73.25)(4,73.31)(5,73.38)};
    \addplot[colorDenseNet,    mark=square]   coordinates{(1,76.81)(2,77.73)(3,77.85)(4,78.03)(5,78.01)};
    \addplot[colorResNet,      mark=triangle, mark options={rotate=180}]
                                              coordinates{(1,77.58)(2,80.41)(3,80.85)(4,81.23)(5,81.34)};
    \addplot[colorEfficientNet,mark=diamond]  coordinates{(1,80.25)(2,80.58)(3,80.66)(4,80.75)(5,80.72)};
    \addplot[colorViT,         mark=*]        coordinates{(1,89.69)(2,90.96)(3,91.31)(4,91.43)(5,91.52)};

    \Ann{south}{1}{66.3}{67.36}{0.87}{colorVGG}
    \Ann{south}{1}{71.6}{72.74}{0.54}{colorMobileNet}
    \Ann{north}{5}{73.8}{73.38}{0.21}{colorMobileNet}
    \Ann{south}{1}{75.7}{76.81}{0.59}{colorDenseNet}
    \Ann{north}{5}{78.4}{78.01}{0.21}{colorDenseNet}
    \Ann{south}{1}{76.4}{77.58}{0.51}{colorResNet}
    \Ann{north}{5}{81.8}{81.34}{0.34}{colorResNet}
    \Ann{north}{1}{80.8}{80.25}{0.32}{colorEfficientNet}
    \Ann{north}{1}{90.2}{89.69}{0.33}{colorViT}
    \Ann{north}{5}{92.0}{91.52}{0.12}{colorViT}

    \SensLegend
  \end{axis}
  \end{tikzpicture}
  \end{adjustbox}
  \caption{Validation Macro-F1 vs.\ $k$}
  \label{fig:tinyimagenet_f1}
\end{subfigure}
\caption{Sensitivity analysis on \textbf{Tiny ImageNet}. The official validation split is used as the held-out evaluation set because test labels are unavailable, so both panels report validation metrics.}
\label{fig:sensitivity_tinyimagenet}
\end{figure}

% ╔══════════════════════════════════════════════════════════════════════╗
% ║  FIGURE 3 — BURMESE GRAPE                                           ║
% ╚══════════════════════════════════════════════════════════════════════╝
\begin{figure}[htbp]
\centering
\begin{subfigure}[t]{0.49\textwidth}
  \begin{adjustbox}{max width=\linewidth}
  \begin{tikzpicture}
  \begin{axis}[sens plot,
    title={Burmese Grape — Accuracy},
    ylabel={Accuracy (\%)},
    ymin=79.5, ymax=90.5]

    \addplot[colorVGG,         mark=x]        coordinates{(1,83.76)(2,85.30)(3,85.43)(4,85.77)(5,85.94)};
    \addplot[colorMobileNet,   mark=triangle] coordinates{(1,83.89)(2,85.94)(3,86.67)(4,87.19)(5,87.01)};
    \addplot[colorDenseNet,    mark=square]   coordinates{(1,84.62)(2,84.96)(3,85.00)(4,84.96)(5,85.17)};
    \addplot[colorResNet,      mark=triangle, mark options={rotate=180}]
                                              coordinates{(1,85.38)(2,85.13)(3,85.17)(4,85.68)(5,85.60)};
    \addplot[colorEfficientNet,mark=diamond]  coordinates{(1,82.74)(2,85.56)(3,85.77)(4,86.54)(5,86.45)};
    \addplot[colorViT,         mark=*]        coordinates{(1,86.98)(2,86.58)(3,87.05)(4,86.79)(5,87.01)};

    \Ann{south}{1}{82.6}{83.76}{2.04}{colorVGG}
    \Ann{north}{5}{86.3}{85.94}{0.98}{colorVGG}
    \Ann{south}{1}{82.7}{83.89}{1.54}{colorMobileNet}
    \Ann{north}{4}{87.6}{87.19}{1.54}{colorMobileNet}
    \Ann{north}{1}{85.0}{84.62}{1.62}{colorDenseNet}
    \Ann{south}{1}{84.3}{85.38}{0.94}{colorResNet}
    \Ann{south}{1}{81.5}{82.74}{2.09}{colorEfficientNet}
    \Ann{north}{4}{87.0}{86.54}{1.05}{colorEfficientNet}
    \Ann{north}{1}{87.5}{86.98}{1.92}{colorViT}

    \SensLegend
  \end{axis}
  \end{tikzpicture}
  \end{adjustbox}
  \caption{Test Accuracy vs.\ $k$}
  \label{fig:grape_acc}
\end{subfigure}\hfill
\begin{subfigure}[t]{0.49\textwidth}
  \begin{adjustbox}{max width=\linewidth}
  \begin{tikzpicture}
  \begin{axis}[sens plot,
    title={Burmese Grape — F1-Score},
    ylabel={F1-Score (\%)},
    ymin=78.0, ymax=90.0]

    \addplot[colorVGG,         mark=x]        coordinates{(1,81.83)(2,83.97)(3,83.69)(4,84.08)(5,84.41)};
    \addplot[colorMobileNet,   mark=triangle] coordinates{(1,81.85)(2,84.55)(3,85.42)(4,85.64)(5,86.08)};
    \addplot[colorDenseNet,    mark=square]   coordinates{(1,82.44)(2,83.01)(3,83.01)(4,82.84)(5,83.06)};
    \addplot[colorResNet,      mark=triangle, mark options={rotate=180}]
                                              coordinates{(1,84.21)(2,84.02)(3,84.08)(4,84.54)(5,84.40)};
    \addplot[colorEfficientNet,mark=diamond]  coordinates{(1,81.17)(2,84.39)(3,84.84)(4,85.49)(5,85.51)};
    \addplot[colorViT,         mark=*]        coordinates{(1,84.85)(2,85.70)(3,86.00)(4,85.66)(5,85.89)};

    \Ann{south}{1}{80.6}{81.83}{3.42}{colorVGG}
    \Ann{north}{5}{84.8}{84.41}{1.42}{colorVGG}
    \Ann{south}{1}{80.6}{81.85}{1.62}{colorMobileNet}
    \Ann{north}{5}{86.5}{86.08}{1.39}{colorMobileNet}
    \Ann{north}{1}{82.9}{82.44}{1.46}{colorDenseNet}
    \Ann{north}{1}{84.7}{84.21}{1.17}{colorResNet}
    \Ann{south}{1}{79.9}{81.17}{2.70}{colorEfficientNet}
    \Ann{north}{5}{85.9}{85.51}{0.90}{colorEfficientNet}
    \Ann{north}{1}{85.3}{84.85}{1.49}{colorViT}
    \Ann{north}{3}{86.5}{86.00}{1.87}{colorViT}

    \SensLegend
  \end{axis}
  \end{tikzpicture}
  \end{adjustbox}
  \caption{Test F1-Score vs.\ $k$}
  \label{fig:grape_f1}
\end{subfigure}
\caption{Sensitivity analysis on \textbf{Burmese Grape}.}
\label{fig:sensitivity_grape}
\end{figure}

% ╔══════════════════════════════════════════════════════════════════════╗
% ║  FIGURE 4 — TOMATO                                                  ║
% ╚══════════════════════════════════════════════════════════════════════╝
\begin{figure}[htbp]
\centering
\begin{subfigure}[t]{0.49\textwidth}
  \begin{adjustbox}{max width=\linewidth}
  \begin{tikzpicture}
  \begin{axis}[sens plot,
    title={Tomato — Accuracy},
    ylabel={Accuracy (\%)},
    ymin=78.5, ymax=91.0]

    \addplot[colorVGG,         mark=x]        coordinates{(1,82.08)(2,83.36)(3,84.00)(4,84.00)(5,83.20)};
    \addplot[colorMobileNet,   mark=triangle] coordinates{(1,83.04)(2,84.40)(3,85.12)(4,85.92)(5,86.64)};
    \addplot[colorDenseNet,    mark=square]   coordinates{(1,84.08)(2,84.56)(3,84.48)(4,84.48)(5,84.40)};
    \addplot[colorResNet,      mark=triangle, mark options={rotate=180}]
                                              coordinates{(1,81.84)(2,85.36)(3,86.64)(4,87.52)(5,87.12)};
    \addplot[colorEfficientNet,mark=diamond]  coordinates{(1,83.20)(2,84.96)(3,85.20)(4,85.52)(5,85.60)};
    \addplot[colorViT,         mark=*]        coordinates{(1,83.12)(2,83.52)(3,83.92)(4,83.68)(5,83.68)};

    \Ann{south}{1}{80.9}{82.08}{2.79}{colorVGG}
    \Ann{north}{3}{84.5}{84.00}{2.56}{colorVGG}
    \Ann{south}{1}{81.8}{83.04}{3.92}{colorMobileNet}
    \Ann{north}{5}{87.1}{86.64}{3.51}{colorMobileNet}
    \Ann{north}{1}{84.6}{84.08}{2.50}{colorDenseNet}
    \Ann{south}{1}{80.6}{81.84}{2.62}{colorResNet}
    \Ann{north}{4}{88.0}{87.52}{2.57}{colorResNet}
    \Ann{north}{1}{83.7}{83.20}{2.46}{colorEfficientNet}
    \Ann{north}{1}{83.7}{83.12}{2.93}{colorViT}

    \SensLegend
  \end{axis}
  \end{tikzpicture}
  \end{adjustbox}
  \caption{Test Accuracy vs.\ $k$}
  \label{fig:tomato_acc}
\end{subfigure}\hfill
\begin{subfigure}[t]{0.49\textwidth}
  \begin{adjustbox}{max width=\linewidth}
  \begin{tikzpicture}
  \begin{axis}[sens plot,
    title={Tomato — F1-Score},
    ylabel={F1-Score (\%)},
    ymin=82.0, ymax=93.5]

    \addplot[colorVGG,         mark=x]        coordinates{(1,84.71)(2,86.04)(3,86.87)(4,86.96)(5,86.43)};
    \addplot[colorMobileNet,   mark=triangle] coordinates{(1,86.78)(2,87.84)(3,88.56)(4,88.93)(5,89.70)};
    \addplot[colorDenseNet,    mark=square]   coordinates{(1,87.52)(2,87.98)(3,87.91)(4,87.94)(5,87.95)};
    \addplot[colorResNet,      mark=triangle, mark options={rotate=180}]
                                              coordinates{(1,84.48)(2,88.08)(3,89.38)(4,90.01)(5,89.91)};
    \addplot[colorEfficientNet,mark=diamond]  coordinates{(1,86.19)(2,88.52)(3,88.58)(4,88.78)(5,88.99)};
    \addplot[colorViT,         mark=*]        coordinates{(1,86.71)(2,87.23)(3,87.37)(4,87.21)(5,87.26)};

    \Ann{south}{1}{83.5}{84.71}{2.67}{colorVGG}
    \Ann{north}{4}{87.4}{86.96}{1.91}{colorVGG}
    \Ann{south}{1}{85.5}{86.78}{3.99}{colorMobileNet}
    \Ann{north}{5}{90.2}{89.70}{2.80}{colorMobileNet}
    \Ann{north}{1}{88.0}{87.52}{1.94}{colorDenseNet}
    \Ann{south}{1}{83.2}{84.48}{1.92}{colorResNet}
    \Ann{north}{4}{90.5}{90.01}{1.84}{colorResNet}
    \Ann{south}{1}{84.9}{86.19}{2.21}{colorEfficientNet}
    \Ann{north}{5}{89.5}{88.99}{1.70}{colorEfficientNet}
    \Ann{north}{1}{87.2}{86.71}{2.77}{colorViT}

    \SensLegend
  \end{axis}
  \end{tikzpicture}
  \end{adjustbox}
  \caption{Test F1-Score vs.\ $k$}
  \label{fig:tomato_f1}
\end{subfigure}
\caption{Sensitivity analysis on \textbf{Tomato}.}
\label{fig:sensitivity_tomato}
\end{figure}

% ╔══════════════════════════════════════════════════════════════════════╗
% ║  FIGURE 5 — POTATO                                                  ║
% ╚══════════════════════════════════════════════════════════════════════╝
\begin{figure}[htbp]
\centering
\begin{subfigure}[t]{0.49\textwidth}
  \begin{adjustbox}{max width=\linewidth}
  \begin{tikzpicture}
  \begin{axis}[sens plot,
    title={Potato — Accuracy},
    ylabel={Accuracy (\%)},
    ymin=83.5, ymax=93.5]

    \addplot[colorVGG,         mark=x]        coordinates{(1,86.03)(2,87.18)(3,88.08)(4,87.61)(5,87.78)};
    \addplot[colorMobileNet,   mark=triangle] coordinates{(1,87.09)(2,88.03)(3,88.29)(4,88.50)(5,88.76)};
    \addplot[colorDenseNet,    mark=square]   coordinates{(1,86.79)(2,89.10)(3,89.66)(4,90.26)(5,90.00)};
    \addplot[colorResNet,      mark=triangle, mark options={rotate=180}]
                                              coordinates{(1,86.20)(2,88.42)(3,88.89)(4,88.76)(5,89.10)};
    \addplot[colorEfficientNet,mark=diamond]  coordinates{(1,85.94)(2,88.12)(3,88.46)(4,88.68)(5,88.97)};
    \addplot[colorViT,         mark=*]        coordinates{(1,90.30)(2,90.30)(3,90.43)(4,90.90)(5,90.60)};

    \Ann{south}{1}{84.8}{86.03}{1.67}{colorVGG}
    \Ann{north}{3}{88.6}{88.08}{1.03}{colorVGG}
    \Ann{south}{1}{85.8}{87.09}{1.67}{colorMobileNet}
    \Ann{north}{5}{89.2}{88.76}{0.72}{colorMobileNet}
    \Ann{south}{1}{85.6}{86.79}{1.88}{colorDenseNet}
    \Ann{north}{4}{90.7}{90.26}{2.01}{colorDenseNet}
    \Ann{south}{1}{85.0}{86.20}{1.02}{colorResNet}
    \Ann{north}{5}{89.6}{89.10}{0.79}{colorResNet}
    \Ann{south}{1}{84.7}{85.94}{0.76}{colorEfficientNet}
    \Ann{north}{5}{89.5}{88.97}{0.69}{colorEfficientNet}
    \Ann{north}{1}{90.8}{90.30}{1.51}{colorViT}
    \Ann{north}{4}{91.4}{90.90}{0.49}{colorViT}

    \SensLegend
  \end{axis}
  \end{tikzpicture}
  \end{adjustbox}
  \caption{Test Accuracy vs.\ $k$}
  \label{fig:potato_acc}
\end{subfigure}\hfill
\begin{subfigure}[t]{0.49\textwidth}
  \begin{adjustbox}{max width=\linewidth}
  \begin{tikzpicture}
  \begin{axis}[sens plot,
    title={Potato — F1-Score},
    ylabel={F1-Score (\%)},
    ymin=82.5, ymax=91.5]

    \addplot[colorVGG,         mark=x]        coordinates{(1,84.52)(2,86.07)(3,87.19)(4,86.68)(5,86.88)};
    \addplot[colorMobileNet,   mark=triangle] coordinates{(1,85.52)(2,86.89)(3,86.92)(4,87.21)(5,87.51)};
    \addplot[colorDenseNet,    mark=square]   coordinates{(1,84.90)(2,87.18)(3,87.64)(4,88.53)(5,88.38)};
    \addplot[colorResNet,      mark=triangle, mark options={rotate=180}]
                                              coordinates{(1,84.84)(2,87.53)(3,87.23)(4,87.50)(5,87.59)};
    \addplot[colorEfficientNet,mark=diamond]  coordinates{(1,84.67)(2,86.80)(3,87.13)(4,87.29)(5,87.57)};
    \addplot[colorViT,         mark=*]        coordinates{(1,88.06)(2,88.29)(3,88.53)(4,89.15)(5,88.94)};

    \Ann{south}{1}{83.3}{84.52}{2.20}{colorVGG}
    \Ann{north}{3}{87.7}{87.19}{1.23}{colorVGG}
    \Ann{south}{1}{84.3}{85.52}{1.94}{colorMobileNet}
    \Ann{north}{5}{88.0}{87.51}{0.81}{colorMobileNet}
    \Ann{south}{1}{83.7}{84.90}{1.90}{colorDenseNet}
    \Ann{north}{4}{89.0}{88.53}{2.22}{colorDenseNet}
    \Ann{south}{1}{83.6}{84.84}{0.68}{colorResNet}
    \Ann{south}{1}{83.4}{84.67}{0.94}{colorEfficientNet}
    \Ann{north}{5}{88.1}{87.57}{1.04}{colorEfficientNet}
    \Ann{north}{1}{88.6}{88.06}{1.31}{colorViT}
    \Ann{north}{4}{89.7}{89.15}{1.21}{colorViT}

    \SensLegend
  \end{axis}
  \end{tikzpicture}
  \end{adjustbox}
  \caption{Test F1-Score vs.\ $k$}
  \label{fig:potato_f1}
\end{subfigure}
\caption{Sensitivity analysis on \textbf{Potato}.}
\label{fig:sensitivity_potato}
\end{figure}


\begin{table}[H]
  \centering
  \caption{Equally weighted mean performance across the 30 model--dataset
    configurations. Each cell reports the aggregate mean in percent, with the
    absolute improvement over the single-checkpoint baseline at $k=1$, in
  percentage points, in parentheses.}
  \label{tab:aggregate_k_results}
  \small
  \setlength{\tabcolsep}{4.5pt}
  \begin{tabular}{lccccc}
    \toprule
    $k$ & Accuracy & Macro-F1 & \makecell{Macro\\Precision}
    & \makecell{Macro\\Recall} & AUC$^{\dagger}$ \\
    \midrule
    1 (Baseline) & 81.73 & 81.65 & 82.58 & 81.58 & 97.77 \\
    2 & 83.10\,{\scriptsize(+1.37)} & 83.19\,{\scriptsize(+1.54)}
    & 83.94\,{\scriptsize(+1.37)} & 83.06\,{\scriptsize(+1.47)}
    & 98.05\,{\scriptsize(+0.28)} \\
    3 & 83.48\,{\scriptsize(+1.76)} & 83.53\,{\scriptsize(+1.88)}
    & 84.17\,{\scriptsize(+1.60)} & 83.44\,{\scriptsize(+1.86)}
    & 98.12\,{\scriptsize(+0.35)} \\
    4 & 83.72\,{\scriptsize(+1.99)} & 83.76\,{\scriptsize(+2.10)}
    & 84.38\,{\scriptsize(+1.80)} & 83.70\,{\scriptsize(+2.11)}
    & 98.15\,{\scriptsize(+0.38)} \\
    5 & \textbf{83.76}\,{\scriptsize(\textbf{+2.04})}
    & \textbf{83.86}\,{\scriptsize(\textbf{+2.21})}
    & \textbf{84.42}\,{\scriptsize(\textbf{+1.85})}
    & \textbf{83.80}\,{\scriptsize(\textbf{+2.22})}
    & \textbf{98.16}\,{\scriptsize(\textbf{+0.39})} \\
    \bottomrule
    \multicolumn{6}{@{}l}{\footnotesize $^{\dagger}$Averaged over the 24
    configurations for which AUC was recorded (all datasets except Tiny}\\
    \multicolumn{6}{@{}l}{\footnotesize \phantom{$^{\dagger}$}ImageNet); the
    other columns average all 30.}
  \end{tabular}
\end{table}

When the configurations are weighted equally, aggregate mean accuracy
increases monotonically from 81.73\% at $k=1$ to 83.76\% at $k=5$,
while aggregate mean macro-F1 increases from 81.65\% to 83.86\%.
These changes correspond to average gains of 2.04 percentage points
in accuracy and 2.21 percentage points in macro-F1. The same monotonic
pattern holds for the remaining metrics: macro-precision gains 1.85
percentage points, macro-recall 2.22, and AUC 0.39. The gain in recall
exceeding the gain in precision indicates that the aggregated model recovers
comparatively more of the under-represented classes, which is consistent with
macro-F1 improving slightly more than accuracy overall. The AUC increase is
small in absolute terms because the baseline already exceeds 97.7\%, yet it is
positive in 23 of the 24 configurations in which AUC is available, indicating
that averaging improves the ranking quality of the predicted scores and not
only the arg-max decision. Thus, although
individual configurations may saturate or decline slightly as $k$
increases, the gains are sufficiently consistent to produce a monotonic
aggregate trend. These results support $k\in\{4,5\}$ as a practical
operating range, although neither value is universally optimal. The
reported aggregates are descriptive macro-averages across heterogeneous
configurations rather than pooled estimates over individual
final-evaluation examples.

All values of $k$ are reported as a sensitivity analysis. The highest
observed results are descriptive and are not used to select $k$ using
the final evaluation sets.

% ── Which checkpoints does the criterion actually select? ────────────────
\FloatBarrier
\subsection{Which Checkpoints Does the Criterion Select?}
\label{sec:selection_geometry}

Sections~\ref{sec:related_work} and~\ref{sec:methodology} distinguish the
proposed criterion from temporal-window averaging on the grounds that a
validation-ranked set need not be contiguous nor confined to the final epochs.
This subsection examines whether that distinction is realised in practice. We
inspect the selected epoch indices for the CIFAR-100 runs, for which the
per-seed selection records are available: six architectures $\times$ five
seeds $=30$ completed runs, and hence 150 selected checkpoints at $k=5$.
Table~\ref{tab:selection_geometry} summarises their positions.

\begin{table}[htbp]
  \centering
  \caption{Position of the selected checkpoints on CIFAR-100, averaged over
    the five seeds. $T$ is the number of epochs actually completed (early
    stopping frequently ends a run before its scheduled budget), $t^{*}$ the
    baseline's single best epoch, and $S_5$ the top-5 validation-ranked set.
    \emph{Span} is $\max S_5-\min S_5+1$, so a purely contiguous selection
    would span exactly 5. \emph{Non-contig.}\ counts seeds whose $S_5$ is not
    a block of consecutive epochs, and the last column counts how many members
  of $S_5$ fall inside the final five epochs of the run.}
  \label{tab:selection_geometry}
  \small
  \setlength{\tabcolsep}{5pt}
  \begin{tabular}{lccccc}
    \toprule
    Model & $T$ & $t^{*}$ & Span of $S_5$
    & \makecell{Non-contig.\\(of 5 seeds)}
    & $|S_5\cap\text{last-5}|$ \\
    \midrule
    VGG16          & 75.8 & 63.8 & 17.6 & 5 & 1.0 \\
    DenseNet121    & 37.2 & 25.2 & 10.0 & 5 & 0.2 \\
    EfficientNetB0 & 23.8 & 11.8 &  9.0 & 5 & 0.2 \\
    MobileNetV2    & 50.0 & 38.0 & 16.0 & 5 & 0.4 \\
    ResNet101      & 46.6 & 34.6 & 14.4 & 5 & 1.6 \\
    ViT-Base       & 14.2 &  6.2 &  5.0 & 0 & 0.0 \\
    \midrule
    All            & 41.3 & 30.0 & 12.0 & 25 & 0.6 \\
    \bottomrule
  \end{tabular}
\end{table}

\paragraph{The selected set is rarely a temporal window.}
A block of five consecutive epochs would have a span of exactly 5. The
observed mean span is 12.0 epochs, and only 5 of the 30 runs---all of them
ViT-Base---produce a contiguous set. In the remaining 25 runs the criterion
skips intervening epochs, with a mean largest internal gap of 5.0 epochs and a
maximum of 10. ViT-Base is the informative exception rather than a
counterexample: it early-stops after roughly 14 epochs, so its five best
checkpoints have little room to be anything other than adjacent. The longer a
run is, the further the selection departs from a window; VGG16, which trains
for 75.8 epochs on average, spreads its five checkpoints over 17.6.

\paragraph{The selected set is not concentrated at the end of training.}
Only 17 of the 150 selected checkpoints ($11.3\%$) lie within the final five
epochs of their run, and in 20 of the 30 runs $S_5$ is entirely disjoint from
that window. On average, the earliest selected checkpoint precedes the end of
training by 16.4 epochs, and the baseline's own best epoch $t^{*}$ sits at
roughly $66\%$ of the completed trajectory rather than at its end---a direct
consequence of combining early stopping with a patience of 8--12 epochs. A
last-$k$ average computed on the same runs would therefore aggregate a set
that is largely disjoint from the one selected here, which is the concrete
sense in which the two criteria differ.

\paragraph{The selection is nested in $k$.}
Because the criterion ranks a fixed set of checkpoints by validation loss,
$S_{k-1}\subset S_k$ holds by construction, and this is confirmed in all 30
runs: increasing $k$ only ever admits an additional checkpoint and never
replaces one. The aggregated model at $k$ is thus a refinement of the model at
$k-1$ rather than a different point in weight space, which explains why the
aggregate curves in Table~\ref{tab:aggregate_k_results} move smoothly rather
than erratically with $k$, even though individual configurations may peak at
an intermediate $k$.

\paragraph{Storage cost.}
The runs retain on average 41.3 checkpoints, since candidates must be kept
until training terminates. This is the method's only material overhead: no
additional optimization is performed, the number of training epochs is
unchanged by construction, and the final model is a single set of averaged
parameters. Retention can be bounded in practice by keeping a running set of
the $k$ best checkpoints observed so far, which requires storing $k+1$ rather
than $T$ states and yields the same result as ranking the full set offline.

\FloatBarrier
\section{Discussion}
\label{sec:discussion}

The matched-trajectory design is central to the interpretation of the
results. For each seed, the single-checkpoint baseline and all Top-$k$
variants are derived from the same completed training run and checkpoint
set. Their differences can therefore be attributed to the post-training
checkpoint-selection rule rather than to changes in architecture,
initialization, optimizer, learning-rate schedule, or training duration.
At the common $k=5$ setting, higher mean accuracy and macro-F1 are observed
in all 30 model--dataset configurations.

A plausible explanation is that validation-guided aggregation reduces
dependence on a single checkpoint whose apparent advantage may partly
reflect temporary variation along the training trajectory. Instead, the
proposed method combines several checkpoints that all achieve comparatively
low validation loss. The consistent gains across convolutional, residual,
densely connected, lightweight, and transformer-based architectures suggest
that the benefit is not restricted to one model family. Because no
additional training or prediction-time ensemble is introduced, the
improvement is obtained entirely at the model-selection stage while
retaining single-model inference cost.

The aggregate and configuration-specific results provide complementary
views of the effect of $k$. As shown in
Table~\ref{tab:aggregate_k_results}, equally weighted mean performance
increases monotonically from $k=1$ to $k=5$, reaching average gains of
2.04 percentage points in accuracy and 2.21 percentage points in macro-F1.
However, individual configurations may saturate or decline slightly after
an intermediate value of $k$. Although $k=4$ or $k=5$ yields the highest
mean accuracy in 25 of the 30 configurations, the remaining five peak at
smaller values. Thus, $k\in\{4,5\}$ represents a useful practical range
rather than a universally optimal choice. If $k$ is tuned for a particular
application, it should be selected using validation data without reference
to the final evaluation results.

Checkpoint aggregation also frequently improves run-to-run stability. At
$k=5$, the accuracy standard deviation is lower than that of the baseline
in 25 of the 30 configurations, and the mean across-seed standard deviation
falls by $22.3\%$ for accuracy and $23.5\%$ for macro-F1. Because several
configurations show small increases, variance reduction should be interpreted
as a general empirical tendency rather than a guaranteed outcome.

The selection analysis in Section~\ref{sec:selection_geometry} clarifies what
the criterion contributes beyond the act of averaging itself. The selected
checkpoints are neither contiguous nor concentrated at the end of training:
only 5 of 30 runs yield a contiguous set, and just $11.3\%$ of the selected
checkpoints fall in the final five epochs. A temporal window of the same size
would therefore aggregate a largely different set of weights on these runs.
This establishes that the two criteria are structurally distinct, but it does
not by itself establish that the validation-ranked set is the better of the
two, since both were not evaluated under a matched protocol here.

The main limitation is that the current experiments establish an advantage
over conventional single-checkpoint selection but do not directly establish
superiority over alternative aggregation strategies such as last-$k$
temporal averaging, EMA, SWA or post-hoc SWA, and prediction ensembling.
Moreover, the evaluation is limited to image classification, so it remains
unclear whether the observed benefits extend to other tasks and modalities,
such as semantic segmentation and natural language processing. The method
also requires candidate checkpoints to be retained until training is
complete, introducing additional storage but no inference-time overhead.
Future work should include controlled comparisons with alternative
aggregation strategies, evaluation on additional tasks and modalities,
adaptive validation-based selection of $k$, diversity-aware aggregation,
and assessment under distribution shifts and input corruptions.

\section{Conclusion}
\label{sec:conclusion}

This study presented a post-training checkpoint aggregation strategy that
ranks epoch checkpoints by validation loss and uniformly averages the
trainable parameters of the top-$k$ candidates. The method operates on a
completed standard training run, requires no secondary optimization phase,
and preserves single-model inference cost.

Across six architectures, five image-classification datasets, and five
random seeds, the common $k=5$ setting achieves higher mean accuracy and
macro-F1 than the single-checkpoint baseline in all 30 model--dataset
configurations. When the configurations are weighted equally, mean accuracy
increases from 81.73\% to 83.76\%, while mean macro-F1 increases from
81.65\% to 83.86\%, corresponding to gains of 2.04 and 2.21 percentage
points, respectively. Checkpoint aggregation also frequently reduces
run-to-run variability.

Although $k=4$ or $k=5$ produces the highest mean accuracy in most
configurations, the effect of the aggregation size is not universally
monotonic. The main empirical finding is therefore not that increasing $k$
always improves performance, but that replacing winner-takes-all checkpoint
selection with a small validation-ranked checkpoint set can provide a more
effective and frequently more stable final model without additional
inference-time complexity. Future work should compare the method directly
with alternative checkpoint-averaging strategies, evaluate its applicability
to other tasks and modalities such as semantic segmentation and natural
language processing, and investigate adaptive or diversity-aware aggregation
under distribution shifts.

\bibliographystyle{splncs04}
\bibliography{Strategy2}

\end{document}

% As shown in the results, validation-loss–guided post-training averaging
% consistently improves classification accuracy across all evaluated architectures.
% On the Tomato dataset, VGG16 achieves an absolute improvement of \textbf{4.0\%}.
% On the Burmese dataset, VGG16 and MobileNetV2 improve by \textbf{3.20\%} and
% \textbf{2.56\%}, respectively, demonstrating the robustness of the proposed
% selection strategy under diverse training configurations.

% \begin{table}[!ht]
% \centering
% \caption{Hyperparameter Configurations for Tomato Classification Models (Validation-loss based)}
% \label{tab:tomato_configs}

% \renewcommand{\arraystretch}{1.3}

% \resizebox{\textwidth}{!}{%
% \begin{tabular}{|l|c|c|c|c|c|}
% \hline
% \textbf{Parameter}
% & \textbf{VGG16}
% & \textbf{ResNet101}
% & \textbf{DenseNet121}
% & \textbf{EfficientNetB0}
% & \textbf{MobileNetV2} \\ \hline

% Batch Size
% & 32
% & 128
% & 32
% & 32
% & 32 \\ \hline

% Num Epochs
% & 200
% & 50
% & 200
% & 200
% & 200 \\ \hline

% Learning Rate
% & $5 \times 10^{-4}$
% & $4 \times 10^{-3}$
% & $5 \times 10^{-4}$
% & $5 \times 10^{-4}$
% & $5 \times 10^{-4}$ \\ \hline

% Weight Decay
% & $5 \times 10^{-4}$
% & $1 \times 10^{-7}$
% & $5 \times 10^{-4}$
% & $5 \times 10^{-4}$
% & $5 \times 10^{-4}$ \\ \hline

% Scheduler
% & \shortstack{ReduceLR\\OnPlateau}
% & \shortstack{Cosine\\Annealing}
% & \shortstack{ReduceLR\\OnPlateau}
% & \shortstack{ReduceLR\\OnPlateau}
% & \shortstack{ReduceLR\\OnPlateau} \\ \hline

% Warmup Strategy
% & X
% & \shortstack{Linear\\Warmup}
% & X
% & X
% & X \\ \hline

% Warmup Epochs
% & X
% & 4
% & X
% & X
% & X \\ \hline

% Early Stopping
% & 15
% & 15
% & 15
% & 15
% & 15 \\ \hline

% LR Decay Patience
% & 5
% & X
% & 5
% & 5
% & 5 \\ \hline

% LR Decay Factor
% & 0.5
% & X
% & 0.5
% & 0.5
% & 0.5 \\ \hline

% ETA Min
% & X
% & $1 \times 10^{-6}$
% & X
% & X
% & X \\ \hline

% Dropout Rate
% & 0.5
% & 0.5
% & 0.5
% & 0.5
% & 0.5 \\ \hline

% Classifier Config
% & [256, 128]
% & [256, 128]
% & [256, 128]
% & [256, 128]
% & [256, 128] \\ \hline
% \end{tabular}%
% }
% \end{table}

% \begin{table}[!ht]
% \centering
% \caption{Hyperparameter Configurations for Burmese Dataset Models (Validation-loss based)}
% \label{tab:burmese_model_configs}

% \renewcommand{\arraystretch}{1.3}

% \resizebox{\textwidth}{!}{%
% \begin{tabular}{|l|c|c|c|c|c|}
% \hline
% \textbf{Parameter}
% & \textbf{VGG16}
% & \textbf{ResNet101}
% & \textbf{MobileNetV2}
% & \textbf{EfficientNetB0}
% & \textbf{DenseNet121} \\ \hline

% Batch Size
% & 32
% & 16
% & 64
% & 32
% & 64 \\ \hline

% Num Epochs
% & 200
% & 200
% & 250
% & 300
% & 250 \\ \hline

% Learning Rate
% & $5 \times 10^{-4}$
% & $4 \times 10^{-4}$
% & $5 \times 10^{-4}$
% & $1 \times 10^{-4}$
% & $5 \times 10^{-4}$ \\ \hline

% Weight Decay
% & $5 \times 10^{-4}$
% & $1 \times 10^{-5}$
% & $5 \times 10^{-4}$
% & $1 \times 10^{-4}$
% & $5 \times 10^{-4}$ \\ \hline

% Scheduler
% & \shortstack{ReduceLR\\OnPlateau}
% & \shortstack{Cosine\\Annealing}
% & \shortstack{ReduceLR\\OnPlateau}
% & \shortstack{ReduceLR\\OnPlateau}
% & \shortstack{ReduceLR\\OnPlateau} \\ \hline

% Warmup Strategy
% & X
% & \shortstack{Linear\\Warmup}
% & X
% & X
% & X \\ \hline

% Warmup Epochs
% & X
% & 10
% & X
% & X
% & X \\ \hline

% Early Stopping
% & 15
% & 20
% & 15
% & 15
% & 15 \\ \hline

% LR Decay Patience
% & 5
% & X
% & 5
% & 5
% & 5 \\ \hline

% LR Decay Factor
% & 0.5
% & X
% & 0.5
% & 0.5
% & 0.5 \\ \hline

% ETA Min
% & X
% & $1 \times 10^{-7}$
% & X
% & X
% & X \\ \hline

% Dropout Rate
% & 0.5
% & 0.4
% & 0.3
% & 0.5
% & 0.3 \\ \hline

% Classifier Config
% & [256, 128]
% & [512, 128]
% & [512, 256]
% & [256, 128]
% & [512, 256] \\ \hline
% \end{tabular}%
% }
% \end{table}
\FloatBarrier
% TABLE 3: TOMATO CONFIGS (TRANSPOSED)
